\begin{enumerate}
\def\labelenumi{(\alph{enumi})}
\setcounter{enumi}{1}
\tightlist
\item
  假设 G 是紧的. 设 \(\phi\) 是 G 的 étale 自同态. 则 \(\operatorname{Ker} \phi\) 是有限的, \(\phi(G)\) 在 G 中的指数有限, 并且
\end{enumerate}

\[
\frac {\operatorname{Card} (\operatorname{Ker} \phi)}{\operatorname{Card} (\mathrm{G} / \phi (\mathrm{G}))} = \operatorname{mod} \det \mathrm{L} (\phi).
\]

(利用 (a) 变换 \(\mu (\mathbf{G})\) .)

\begin{enumerate}
\def\labelenumi{(\alph{enumi})}
\setcounter{enumi}{2}
\tightlist
\item
  假设 G 是紧的且交换的. 设 \(r \in \mathbf{Z}\), 使得 r, 1 在 K 中均不为 0, 并设 \(r, 1 \neq 0\)\(\phi\) 是 G 上的自同态 \(x \mapsto x^r\). 则
\end{enumerate}

\[
\frac {\operatorname{Card} (\operatorname{Ker} \phi)}{\operatorname{Card} (\mathrm{G} / \phi (\mathrm{G}))} = (\mathrm{mod} r. 1) ^ {n}.
\]

(注意由 §2 命题 7 的推论, \(\mathrm{L}(\phi)\) 是比为 r 的伸缩变换.)\(r\)

\begin{enumerate}
\def\labelenumi{\arabic{enumi}.}
\setcounter{enumi}{2}
\tightlist
\item
  设 E 是由 2 行 2 列复对称矩阵组成的复向量空间. 对所有 \(s \in \mathbf{SL}(2, \mathbf{C})\), 定义 E 的自同构 \(\varphi(s)\), 使 \(\varphi(s).m = s.m.^t s\) .
\end{enumerate}

\begin{enumerate}
\def\labelenumi{(\alph{enumi})}
\tightlist
\item
  证明 \(\rho\) 是 \(\mathbf{SL}(2, \mathbf{C})\) 的解析线性表示. 将 \(\begin{pmatrix} a & b \\ b & c \end{pmatrix} \in \mathrm{E}\) 识别为 \((a, b, c) \in \mathbf{C}^3\) 后, \(\rho\begin{pmatrix} \alpha & \beta \\ \gamma & \delta \end{pmatrix}\) 的矩阵为
\end{enumerate}

\[
\left( \begin{array}{c c c} \alpha^ {2} & 2 \alpha \beta & \beta^ {2} \\ \alpha \gamma & \alpha \delta + \beta \gamma & \beta \delta \\ \gamma^ {2} & 2 \gamma \delta & \delta^ {2} \end{array} \right).
\]

\begin{enumerate}
\def\labelenumi{(\alph{enumi})}
\setcounter{enumi}{1}
\tightlist
\item
  证明 \(m \mapsto \det(m)\) 是非退化二次型
\end{enumerate}

\[
q (a, b, c) = a c - b ^ {2}
\]

在 E 上, 并且 \(\rho\) 是从 \(\mathbf{SL}(2,\mathbf{C})\) 到 \(\mathbf{SO}(q)\) 的满射态射, 核为 \(\{I, - I\}\) . (关于满射性, 注意 \(\mathbf{SO}(q)\) 中一个元素的矩阵的第 1 列可表示为 \((\alpha^2,\alpha \gamma ,\gamma^2)\) 的形式, 第 3 列可表示为 \((\beta^2,\beta \delta ,\delta^2)\) 的形式, 其中 \(\alpha^2\delta^2 +\beta^2\gamma^2 -2\alpha \gamma \beta \delta = 1\) . 然后注意, 在一个非各向同性平面上相同的 \(\mathbf{SO}(q)\) 中两个元素彼此相等.)

\begin{enumerate}
\def\labelenumi{(\alph{enumi})}
\setcounter{enumi}{2}
\tightlist
\item
  推出复李群 \(\mathbf{SO}(3, \mathbf{C})\) 同构于复李群 \(\mathbf{SL}(2, \mathbf{C}) / \{I, -I\}\) .
\end{enumerate}

\begin{enumerate}
\def\labelenumi{\arabic{enumi}.}
\setcounter{enumi}{3}
\tightlist
\item
  \begin{enumerate}
  \def\labelenumii{(\alph{enumii})}
  \tightlist
  \item
    设 \(q\) 是 \(\mathbf{R}^3\) 上符号为 (1, 2) 的二次型. 令 \(\mathbf{P}\) 为满足  的 \(m \in \mathbf{R}^3\) 的集合. 则 \(q(m) > 0\)\(\mathbf{P}\) 是两个不交凸锥 \(\mathbf{C}\) 和 \(-\mathbf{C}\) 的并. 若 \(s \in \mathbf{SO}(q)\), 则或者 \(s(\mathbf{C}) = \mathbf{C}\) 且 \(s(-\mathbf{C}) = -\mathbf{C}\), 或者 \(s(\mathbf{C}) = -\mathbf{C}\) 且 \(s(-\mathbf{C}) = \mathbf{C}\). 令 \(\mathbf{SO}^{+}(q)\) 为满足  的 \(s \in \mathbf{SO}(q)\) 的集合. 则 \(s(\mathbf{C}) = \mathbf{C}\)\(\mathbf{SO}^{+}(q)\) 是 \(\mathbf{SO}(q)\) 的开正规子群, 且在 \(\mathbf{SO}(q)\) 中的指数为 2 .
  \end{enumerate}
\end{enumerate}

\begin{enumerate}
\def\labelenumi{(\alph{enumi})}
\setcounter{enumi}{1}
\item
  仿照练习 3 的方法, 定义一个从 \(\mathbf{SL}(2,\mathbf{R})\) 到 \(\mathbf{SO}^{+}(q)\) 的满射态射, 其核为 \(\{I, - I\}\) . 推出实李群 \(\mathbf{SO}^{+}(q)\) 同构于实李群 \(\mathbf{SL}(2,\mathbf{R}) / \{I, - I\}\) .
\item
  对于 \((\xi, \eta) \in \mathbf{C}^2\) 和 \((\xi', \eta') \in \mathbf{C}^2\), 令 \(f((\xi, \eta), (\xi', \eta')) = \xi \bar{\xi}' - \eta \bar{\eta}'\) . 证明由  定义的 \(\mathbf{SL}(2, \mathbf{C})\) 的内自同构将 \(\frac{1}{\sqrt{2}} \begin{pmatrix} 1 & -i \\ -i & 1 \end{pmatrix}\)\(\mathbf{SU}(f)\) 映到 \(\mathbf{SL}(2, \mathbf{R})\) .
\end{enumerate}

\begin{enumerate}
\def\labelenumi{\arabic{enumi}.}
\setcounter{enumi}{4}
\tightlist
\item
  设 E 是由迹为零的 2 行 2 列厄米复矩阵组成的实向量空间. 对所有 \(s \in \mathbf{SU}(2, \mathbf{C})\), 定义 E 的自同构 \(\rho(s)\), 使 \(\rho(s)(m) = sms^*\) . 将元素
\end{enumerate}

\[
\left( \begin{array}{c c} x & y + i z \\ y - i z & - x \end{array} \right)
\]

与  中的元素 \((x,y,z)\) 对应. 仿照练习 3 的方法, 证明 \(R^{3}\)\(\rho\) 是从实李群 \(\mathbf{SU}(2,\mathbf{C})\) 到实李群 \(\mathbf{SO}(3,\mathbf{R})\) 的满射态射, 核为 \(\{I,-I\}\), 并且 \(\mathbf{SO}(3,\mathbf{R})\) 同构于 \(\mathbf{SU}(2,\mathbf{C})/\{I,-I\}\)

¶ 6. 设 F 是由 2 行 2 列厄米复矩阵组成的向量空间. 对所有 \(s \in \mathbf{SL}(2, \mathbf{C})\), 定义 F 的自同构 \(\sigma(s)\), 使 \(\sigma(s)(m) = sms^*\) . 将 F 中的元素 \(\begin{pmatrix} t + x & y + iz \\ y - iz & t - x \end{pmatrix}\) 与  中的元素 \((t, x, y, z)\) 对应, 并记\(\mathbf{R}^4\)

\[
q (t, x, y, z) = t ^ {2} - x ^ {2} - y ^ {2} - z ^ {2}.
\]

令 C 为满足 \((t, x, y, z) \in \mathbf{R}^+\)\(t^2 - x^2 - y^2 - z^2 > 0\), \(t > 0\) 的  的集合. 令 \(\mathbf{SO}^+(q)\) 为满足  的 \(g \in \mathbf{SO}(q)\) 的集合; 它是 \(g(\mathbf{C}) = \mathbf{C}\)\(\mathbf{SO}(q)\) 的开正规子群, 在 \(\mathbf{SO}(q)\) 中的指数为 2 (参见练习 4). 证明 \(\sigma\) 是从 \(\mathbf{SL}(2, \mathbf{C})\) 的底层实李群到实李群 \(\mathbf{SO}^+(q)\) 的满射态射, 核为 \(\{I, -I\}\) . (关于满射性, 使用练习 5.) 推出实李群 \(\mathbf{SO}^+(q)\) 同构于 \(\mathbf{SL}(2, \mathbf{C}) / \{I, -I\}\) 的底层实李群, 即 (练习 3) \(\mathbf{SO}(3, \mathbf{C})\) 的底层实李群 .

\begin{enumerate}
\def\labelenumi{\arabic{enumi}.}
\setcounter{enumi}{6}
\tightlist
\item
  \begin{enumerate}
  \def\labelenumii{(\alph{enumii})}
  \tightlist
  \item
    令 G 为范数为 1 的四元数集合. 这是实李群 \(\mathbf{H}^*\) 的李子群, 同胚于 \(\mathbf{S}_3\), 因而单连通, 参见《一般拓扑》第 XI 章.
  \end{enumerate}
\end{enumerate}

\begin{enumerate}
\def\labelenumi{(\alph{enumi})}
\setcounter{enumi}{1}
\tightlist
\item
  设 E 是纯四元数组成的实向量空间, 通过下式与 \(\mathbf{R}^3\) 对应
\end{enumerate}

\[
(x, y, z) \mapsto x i + y j + z k.
\]

对所有 \(g \in \mathbf{G}\), 定义  的自同构 \(\rho(g)\), 使 \(\mathbf{E}\)\(\rho(g)q = ggg^{-1}\) . 证明 \(\rho\) 是从李群 \(\mathbf{G}\) 到实李群 \(\mathbf{SO}(3, \mathbf{R})\) 的满射态射, 核为 \(\{1, -1\}\) . 推出李群 \(\mathbf{SO}(3, \mathbf{R})\) 同构于李群 \(G / \{1, -1\}\) .

\begin{enumerate}
\def\labelenumi{(\alph{enumi})}
\setcounter{enumi}{2}
\item
  将 \(\mathbf{C}\) 与  的子域 \(\mathbf{R} + \mathbf{R}i\) 对应. 映射 \(\mathbf{H}\)\((\lambda, q) \mapsto q\lambda\) 将 \(\mathbf{C} \times \mathbf{H}\) 映到 \(\mathbf{H}\), 使 \(\mathbf{H}\) 成为复向量空间; 通过选取基 \(\mathbf{C}\), 它与 \(\mathbf{C}^2\) 对应. 对所有 \((1, j)\)\(g \in \mathbf{G}\), 令 \(\sigma(g)\) 为由 \(\sigma(g)q = gq\) 定义的该向量空间的自同构. 证明 \(\sigma\) 是从实李群 G 到实李群 \(\mathbf{SU}(2, \mathbf{C})\) 的同构. 因而由 \((a)\), 后者是单连通的.
\item
  对于 \((q_{1}, q_{2}) \in \mathbf{G} \times \mathbf{G}\), 令 \(\tau(q_{1}, q_{2})\) 为从  到自身的映射 \(q \mapsto q_{1}q\bar{q}_{2}\). 证明 \(\mathbf{H} = \mathbf{R}^{4}\)\(\tau\) 是从实李群 \(\mathrm{G} \times \mathrm{G}\) 到实李群 \(\mathbf{SO}(4, \mathbf{R})\) 的满射态射, 核为 \(\mathrm{N} = \{(1, 1), (-1, -1)\}\), 因而 \(\mathbf{SO}(4, \mathbf{R})\) 同构于 \((\mathrm{G} \times \mathrm{G}) / \mathrm{N}\) .
\end{enumerate}

\begin{enumerate}
\def\labelenumi{\arabic{enumi}.}
\setcounter{enumi}{7}
\tightlist
\item
  设 \(\mathrm{G}\) 是有限维连通实李群, \(\mathbf{M}\) 是有限维连通的 \(\mathrm{C}^{\infty}\) 类流形; 给定 \(\mathrm{C}^{\infty}\)\(\mathrm{G}\) 在 \(\mathbf{M}\) 上的  类右作用法则. 对所有 \(x \in \mathbf{M}\), 令 \(\varphi(x)\) 为相应的轨道映射.
\end{enumerate}

\begin{enumerate}
\def\labelenumi{(\alph{enumi})}
\item
  G 在 M 上传递作用, 当且仅当对所有 \(x \in M\), 从 L(G) 到  的映射 \(\mathrm{T}_{e}(\varphi(x))\) 是满射. (要看条件的充分性, 注意此时 M 中的每个 G-轨道都是开集.)\(\mathrm{T}_{x}(\mathrm{M})\)
\item
  假设 G 在 M 上传递作用. 我们把 M 识别为齐性空间 H\textbar G, 其中 H 是 M 中一点 \(x_{0}\) 的稳定子. 若存在 M 的一个 \(C^{\omega}\) 类闭子流形 V, 满足 \(0 < dim V < dim M\), 且对每个变换 V.s (其中 \(s \in G\)) 而言, V.s 要么等于 V, 要么与 V 不交, 则称 G 在 M 上的作用是非本原的. 这当且仅当存在 G 的李子群 L, 使得 \(H \subset L \subset G\) 且 \(\dim(H) < \dim(L) < \dim(G)\). 若不存在这样的 L, 则称 G 在 M 上的作用是本原的; 只要 L(G) 与 L(H) 之间不存在不同于 L(G) 和 L(H) 的子代数, 这就足够了.
\item
  在 (b) 的假设下, 令 \(\mathcal{L}\) 为 L(G) 在与 G 的作用相联系的无穷小作用法则下的像. 令 \(\mathcal{E}\) 为 M 上 \(C^\infty\) 类函数的代数, m 为 \(\mathcal{E}\) 的极大理想, 由 \(\mathcal{E}\) 中在 \(x_0\) 处为零的函数组成; 对于 , 记 \(\mathcal{L}_p\) 为满足对所有 \(p = -1, 0, 1, 2, \ldots\) 都有  的 \(x \in \mathcal{L}\) 的集合; 于是 \(\theta(\mathbf{X}) f \in m^{p+1}\)\(f \in \mathcal{E}\)\(\mathcal{L}_{-1} = \mathcal{L}\), 而 \(\mathcal{L}_0\) 是 L(H) 在无穷小作用法则下的像. 若 (U, \(\phi, \mathbf{R}^n\) ) 是以 \(x_0\) 为中心的 M 的图, 且 \(\mathbf{X}_1, \ldots, \mathbf{X}_n\) 是由该图定义的 U 上向量场, 则 \(\mathcal{L}_p\) 的元素是形如 \(a_1\mathbf{X}_1 + \cdots + a_n\mathbf{X}_n\) 的元素, 其中 \(a_1, \ldots, a_n\) 属于 \(m^p|\text{U}\). 证明 \([\mathcal{L}_p, \mathcal{L}_q] \subset \mathcal{L}_{p,q}\) (约定 \(\mathcal{L}_{-2} = \mathcal{L}\)). 若存在 Y \(\in \mathcal{L}_p\) (其中 \(p \geqslant 0\)) 使得 Y \(\notin \mathcal{L}_{p+1}\), 则存在 X \(\in\) L 使得 {[}Y, X{]} \(\notin \mathcal{L}_p\). \(\mathcal{L}_p\) 是 \(\mathcal{L}\) 的李子代数. 对于 \(p \geqslant 0\), \(\mathcal{L}_p\) 是 \(\mathcal{L}_0\) 的理想. 令 \(\rho\) 为 H 在 T \(_{x_0}\) (M) 上的典范线性表示. H/(Ker \(\rho\) ) 的李代数同构于 \(\mathcal{L}_0/\mathcal{L}_1\).
\item
  进一步假设各 \(\mathcal{L}_p\) 的交为 \(\{0\}\). 则存在最大的指标 r, 使 \(r\)\(\mathcal{L}_r \neq \{0\}\), 且所有指标  的 \(\mathcal{L}_f\) 两两不同. 对于 \(\leqslant r\)\(p \geqslant 0\),
\end{enumerate}

\[
\dim (\mathcal {L} _ {p} / \mathcal {L} _ {p + 1}) \leqslant \binom {n + p} {p + 1}.
\]

\begin{enumerate}
\def\labelenumi{(\alph{enumi})}
\setcounter{enumi}{4}
\tightlist
\item
  若 M 是解析流形且 G 在 M 上解析地作用, 则满足 (d) 的假设.
\end{enumerate}

\begin{enumerate}
\def\labelenumi{\arabic{enumi}.}
\setcounter{enumi}{8}
\tightlist
\item
  在命题 30 的记号下, HH' 可以是 G 中稠密而不同于 G 的开子集 (取 G = GL(2, R), H 为上三角子群, H' 为严格下三角子群).
\end{enumerate}

\exercisesfor{4}

\begin{enumerate}
\def\labelenumi{\arabic{enumi}.}
\tightlist
\item
  设 G 是一个李群, H 是 G 的正规李拟子群, \(\pi\) 是从 G 到 G/H 的典范映射. 在 G/H 上存在唯一的李群结构, 具有如下性质: 对于从 G/H 到李群 G' 的同态 \(\theta\), \(G'\) 是李群态射当且仅当 \(\theta \circ \pi\) 是李群态射. 此外, L(G/H) 典范地同构于 L(G)/L(H). (令 Q 为满足 L(Q) = L(G)/L(H) 的李群芽. 证明必要时缩小 Q 后, Q 可识别为 G/H 中 0 的一个开邻域, 从而 §1 的命题 18 可应用于 G/H.)
\end{enumerate}

¶ 2. 设 G 和 H 是两个李群, f 是从 G 到 H 的李群态射.

\begin{enumerate}
\def\labelenumi{(\roman{enumi})}
\item
   的核 \(\mathbf{N}\) 是 \(f\)\(\mathbf{G}\) 的正规李拟子群, 且 \(\mathbf{L}(\mathbf{N}) = \operatorname{Ker} \mathbf{L}(f)\) . (使用 \(\mathbf{G}\) 和 \(\mathbf{H}\) 的指数映射.)
\item
  令 g \(g \colon \mathrm{G} / \mathrm{N} \to f(\mathrm{G})\) 为 f 取商后导出的映射. 若 f 和 \(f\)\(f\)\(\mathrm{L}(f)\) 的像都是闭的, 且 \(\mathrm{G}\) 的拓扑有可数基, 则 f(\(f(\mathrm{G})\)G) 是 \(\mathrm{H}\) 的李拟子群, 其李代数为 \(\operatorname{Im} \mathrm{L}(f)\), 且 g 是李群同构, 其中 \(g\)\(\mathrm{G} / \mathrm{N}\) 具有练习 1 中定义的结构. (利用 (i) 将问题归约为 \(\mathrm{N} = \{e\}\) 的情形.)
\end{enumerate}

\begin{enumerate}
\def\labelenumi{\arabic{enumi}.}
\setcounter{enumi}{2}
\tightlist
\item
  设 G 是一个李群, U 是 L(G) 中 0 的开邻域, \(\phi\) 是从 U 到 G 的解析映射, 满足 \(\phi(0) = e\) 且 \(T_{0}\phi = \mathrm{Id}_{\mathrm{L}(G)}\). 以下条件等价:
\end{enumerate}

\begin{enumerate}
\def\labelenumi{(\alph{enumi})}
\item
  \(\phi\) 是指数映射;
\item
  存在不同于 0, 1, -1 的 \(m \in \mathbf{Z}\), 使得在 0 的某个邻域内 \(\phi(mx) = \phi(x)^m\).
\end{enumerate}

(若条件 (b) 成立, 证明 \(\phi^{*}\) 满足定理 4 的条件 (iii).)

¶ 4. (a) 设 K = R 或 C, 或剩余域特征为 \textgreater0 的超度量域. 设 G 是 K 上的李群, U 是 L(G) 中 0 的开邻域, \(\phi: U \to G\) 是在 0 处可微的映射, 并且
满足 \(\phi(0) = e\), \(T_0(\phi) = \mathrm{Id}_{L(G)}\), 且当 \(\phi((\lambda + \lambda')b) = \phi(\lambda b)\phi(\lambda'b)\)\(\lambda b\), \(\lambda'b\), \((\lambda + \lambda')b\) 属于 U 时有 . 则 \(\phi\) 在 0 的某个邻域上与一个指数映射重合.

\begin{enumerate}
\def\labelenumi{(\alph{enumi})}
\setcounter{enumi}{1}
\tightlist
\item
  设 k 是特征为 0 的域. 令 \(k\)\(\mathbf{K}\) 为赋值域 \(k((\mathbf{X}))\), 它的特征为 0. 令 \(\mathbf{G}\) 为加法李群 \(k[[\mathbf{X}]]\). 设 \(\phi\) 是从 \(k\)\(\mathbf{G}\) 到 \(\mathbf{G}\) 的连续 k-线性映射, 满足对每个整数  都有 \(\phi(\mathbf{X}^n) = \mathbf{X}^n + \mathbf{X}^{2n}\). 则 \(n \geqslant 0\)\(\phi\) 满足 (a) 的条件, 但在 0 的任何邻域上都不与指数映射重合.
\end{enumerate}

\begin{enumerate}
\def\labelenumi{\arabic{enumi}.}
\setcounter{enumi}{4}
\tightlist
\item
  设 \((e_1, e_2, e_3)\) 是 \(\mathbf{R}^3\) 的典范基. 在 \(\mathbf{R}^3\) 上赋予幂零李代数结构 \(\mathfrak{g}\), 使得 \([e_1, e_2] = e_3\), \([e_1, e_3] = [e_2, e_3] = 0\). 令 \(c_{ijk}\) 表示相对于基  的 \(\mathfrak{g}\) 的结构常数.\((e_1, e_2, e_3)\)
\end{enumerate}

\begin{enumerate}
\def\labelenumi{(\alph{enumi})}
\tightlist
\item
  证明在 \(\mathbf{R}^3\) 上微分形式
\end{enumerate}

\[
\omega_ {1} = d x _ {1}, \quad \omega_ {2} = d x _ {2}, \quad \omega_ {3} = - x _ {1} d x _ {2} + d x _ {3}
\]

满足关系

\[
d \omega_ {k} = - \sum_ {i <   j} c _ {i j k} \omega_ {i} \wedge \omega_ {j} \quad (k = 1, 2, 3)
\]

并且在 \(\mathbf{R}^3\) 的每一点线性无关.

\begin{enumerate}
\def\labelenumi{(\alph{enumi})}
\setcounter{enumi}{1}
\tightlist
\item
  推出在 \(\mathbf{R}^3\) 中 0 的某个开邻域上存在李群芽结构 \(\mathbf{G}\), 使得 \(\mathrm{L}(\mathrm{G}) = \mathfrak{g}\), 且
\end{enumerate}

\[
(a _ {1}, a _ {2}, a _ {3}) (x _ {1}, x _ {2}, x _ {3}) = (x _ {1} + a _ {1}, x _ {2} + a _ {2}, x _ {3} + a _ {3} + a _ {1} x _ {2})\tag{1}
\]

当 \((a_{1}, a_{2}, a_{3})\), \((x_{1}, x_{2}, x_{3})\) 充分接近 0 时成立. 证明事实上公式 (1) 在 \(\mathbf{R}^{3}\) 上定义了幂零李群结构.

¶ 6. 设 G 是有限维李群, g 是其李代数, g* 是 g 的对偶向量空间, σ 是 G 在 g 上的伴随表示, ρ 是 G 在 g* 上的表示, 对所有 g ∈ G 定义为 ρ(g) = tσ(g)-1.

\begin{enumerate}
\def\labelenumi{(\alph{enumi})}
\item
  对所有 , \(\mathrm{L}(\rho)(x) = -^{t}(\mathrm{ad} x)\).\(x \in g\)
\item
  设 \(f \in \mathfrak{g}^*\). 令 \(G_f\) 表示 f 在 G 中的稳定子; 它是 G 的李子群, 且 \(f\)\(G\)\(G\)\(\mathfrak{g}_f = L(G_f)\) 是满足  的 \(x \in G\) 的集合.\({}^t (\operatorname{ad}x)f = 0\)
\item
  对 \(x, y\), 记 \(\mathfrak{g}\)\(\mathrm{B}_f(x, y) = f([x, y])\). 证明 \(\mathrm{B}_f\) 是 \(\mathfrak{g}\) 上的交错双线性形式, 与 \(s_f\) 左相关的从 \(\mathfrak{g}\) 到 \(\mathfrak{g}^*\) 的映射  是映射  , 且相对于 \(\mathrm{B}_f\), \(x \mapsto {}^t (\mathrm{ad}x)\)\(f\)\(\mathfrak{g}\) 的正交补是 \(\mathrm{B}_f\)\(\mathfrak{g}_f\). 令 \(\beta_f\) 表示在取商时由  导出的 \(\mathfrak{g} / \mathfrak{g}_f\) 上的非退化交错双线性形式 (因而其维数为偶数). 我们将  记为  上 \(\mathrm{B}_f\) 的逆形式.\(\beta_f'\)\(\mathrm{B}_f\)\((\mathfrak{g} / \mathfrak{g}_f)*\)
\item
  令 \(\Omega\) 表示 G 在 \(G\)\(\mathfrak{g}^*\) 上的一个轨道; 通过结构传递, 赋予它由 \(G / G_f\) 上的结构导出的流形结构, 其中 f 是 \(f\)\(\Omega\) 的任意一点. 然后 G 在 \(G\)\(\Omega\) 上解析地左作用. 令 D 为相应的无穷小作用法则. 对所有 \(D\)\(a \in \mathfrak{g}\) 和所有 \(f \in \Omega\),
\end{enumerate}

\[
\mathrm{D} _ {a} (f) = - (^ {t} \mathrm{ad} a) f.
\]

切子空间 \(\mathrm{T}_{f}(\Omega)\) 是 \(s_{f}\) 的像, 即  中 \(g_{f}\) 的正交补, 并且典范地与 \(g^{*}\)\((\mathfrak{g}/\mathfrak{g}_{f})^{*}\) 对应, 因而 \(s_{f}\) 定义了从  到其对偶的典范同构 \(r_{f}\). 场 \(\mathrm{T}_{f}(\Omega)\)\(f \mapsto \beta_{f}'\) 是  上的 2 次解析微分形式 \(\omega\), 且在 G 下不变. 对于 \(\Omega\)\(f \in \Omega\), \(a \in g\), \(b \in g\), \(\omega(\mathrm{D}_{a}(f), \mathrm{D}_{b}(f)) = f([a, b])\).

\begin{enumerate}
\def\labelenumi{(\alph{enumi})}
\setcounter{enumi}{4}
\item
  证明 \(d\omega = 0\) .
\item
  若 \(\alpha\) 是 \(\Omega\) 上的 1 次微分形式, 同构 \(r_f\) 使我们能够将 \(\alpha\) 识别为一个向量场. 令 A 为取值于 K 的 \(\Omega\) 上解析函数的集合. 对于 A 中的 \(\phi, \psi\), 记 \([\phi, \psi] = \omega(d\phi, d\psi) \in A\) . 证明如此定义的 A 具有李代数结构.
\item
  对所有 \(a \in \mathfrak{g}\), 令 \(\psi_{a}\) 为  上的函数 \(f \mapsto f(a)\). 证明 \(\Omega\)\(a \mapsto \psi_{a}\) 是从 \(\mathfrak{g}\) 到李代数 A 的同态, 且形式 \(d\psi_{a}\) 通过同构 \(r_{f}\) 与向量场 \(\mathrm{D}_{a}\) 对应.
\item
  令 U 是 \(\mathfrak{g}^*\) 的开子集, \(\phi\) 是 U 上的解析函数. 对于所有 \(f\in \mathrm{U}\), \(d_{f}\phi\)\(\phi\) 在 \(f\) 处的微分  被识别为 \(\mathfrak{g}\) 的一个元素. 假设对 G 在  上的作用所定义的每个向量场 X 都有 \(\mathrm{X}\phi = 0\). 证明于是对所有 \(\mathfrak{g}^*\)\(f\in \mathrm{U}\), \(d_{f}\phi\) 属于 \(\mathfrak{g}_f\) 的中心.
\item
  对所有 \(f \in \mathfrak{g}^*\), 记 \(r_f = \dim \mathfrak{g}_f\). 令 \(r = \inf_{f \in \mathfrak{g}^*} r_f\). 证明集合 V = \(V\) 在 \(f \in \mathfrak{g}^*\)\(r_f = r\)\(\mathfrak{g}^*\) 中开且稠密. 证明对所有 \(f \in V\), \(\mathfrak{g}_f\) 是交换的. (在 f 的某个邻域内构造 r 个满足 (h) 条件的函数 \(r\)\(\phi\), 并使它们在 f 处的微分线性无关.)\((h)\)\(f\)\(f\)
\end{enumerate}

¶ 7. (a) 设 \((\lambda, x) \mapsto \lambda.x\) 是 R 在 T 上的连续左作用法则. 则以下两种情形之一成立:

\begin{enumerate}
\def\labelenumi{(\arabic{enumi})}
\item
  要么存在一个不动点, 且每一点的稳定子或者是 \(\mathbf{R}\), 或者是 \(\{0\}\) ;
\item
  要么存在一点, 其稳定子是 \(\mathbf{R}\) 的无限离散子群, 且 \(\mathbf{R}\) 在 \(\mathbf{T}\) 上传递作用. (若 \(\{\lambda \in \mathbf{R}|\lambda .x = x\} = \{0\}\), 则 x 的轨道同胚于 \(x\)\(\mathbf{R}\), 并具有边界点 \(\lim_{\lambda \to +\infty}\lambda .x\), 这些边界点是不动点. 若 \(\{\lambda \in \mathbf{R}|\lambda .x = x\} = \mathbf{Z}a\) 且 \(a\neq 0\), 则 x 的轨道同胚于 \(x\)\(\mathbf{T}\), 且 \(\mathbf{R}\) 是传递的.)
\end{enumerate}

\begin{enumerate}
\def\labelenumi{(\alph{enumi})}
\setcounter{enumi}{1}
\item
  若 f 是 \(f\)\(\mathbf{T}\) 到自身的同胚, k 是一个 \(k\)\(\geqslant 1\) 的整数, 则称 \(x \in \mathbf{T}\) 是 f 的周期为 k 的周期点, 如果 \(k\)\(f\)\(f^k(x) = x\) 且对  有 \(f^h(x) \neq x\). 在 (a) 的记号下, 令 \(1 \leqslant h < k\)\(f_{\lambda}\) 为映射 \(x \mapsto \lambda .x\) 所定义的同胚. 则 \(k > 1\)\(f_{\lambda}\) 的周期大于 1 且等于 k 的周期点集为空或等于 \(\mathbf{T}\). (若  存在周期为 \(k > 1\) 的周期点, 则其稳定子不同于 \(f_{\lambda}\)\{0\} 和 \(\mathbf{R}\), 因而 \(\mathbf{R}\) 在 \(\mathbf{T}\) 上传递作用, 且 \(\mathbf{T}\) 的所有点都是 \(k\) 的周期为 k 的周期点.)\(f_{\lambda}.\)
\item
  令 \(\Omega\) 是 \(e\)\(\mathrm{Diff}^{\infty}(\mathbf{T})\) 中 e 的一个邻域 (《可微与解析流形》, R, 15.3.8, 注 (1)). 存在整数 \(k > 1\) 以及无不动点的 \(f\in \Omega\), 使得 f 的周期为 k 的周期点集既非空又不等于 \(k\)\(f\)\(\mathbf{T}\). (令 \(p\colon \mathbf{R}\to \mathbf{T}\) 为典范映射. 令 \(\rho\) 为  上的映射 \(p(x) \mapsto p\left(x + \frac{2\pi}{k}\right)\). 令 \(\mathbf{T}\)\(\mathbf{T}\)\(\Omega'\) 是 \(e\)\(\text{Diff}^\infty(\mathbf{T})\) 中 e 的一个邻域, 使 \(\Omega'^2 \subset \Omega\). 当 k 充分大时, \(k, \rho \in \Omega'\). 另一方面, 令 \(\sigma\) 是 \(\mathbf{T}\) 到自身的映射, 使得 \(\mathbf{T}\)\(\sigma(y) = y\) 对于
\end{enumerate}

\[
y \notin p \left(\right] 0, \frac {2 \pi}{k} [)
\]

且当  时 \(\sigma(p(x)) = p(x + \lambda(x))\), 其中\(x \in \left]0, \frac{2\pi}{k}\right\}\)

\[
0 <   \lambda (x) <   \frac {2 \pi}{k} - x;
\]

可以选取 \(\sigma \in \Omega'\). 则 \(f = \rho \circ \sigma\) 具有所要求的性质.) 根据 (b), 这样的 f 不可能属于从 \(f\)\(\mathbf{R}\) 到 \(\operatorname{Diff}^{\infty}(\mathbf{T})\) 的连续态射的像.

\begin{enumerate}
\def\labelenumi{(\alph{enumi})}
\setcounter{enumi}{3}
\item
  设 Y 是维数 \(\geqslant1\) 的紧可微流形. \(\operatorname{Diff}^{\infty}(Y)\) 的每个 e 邻域都包含一个具有如下性质的元素 h: h 不属于从 R 到 \(\operatorname{Diff}^{0}(Y)\) 的任何连续态射的像. (流形 Y 包含一个形如 \(T \times D\) 的开子流形, 其中 D 是 \(R^{n}\) 中半径为 1 且中心为 0 的开欧氏球 ( \(n = \dim Y - 1\) ). 取 \(g \in \operatorname{Diff}^{\infty}(D)\), 使得 \(g(0) = 0\), 当  时 \(\|g(x)\| > \|x\|\), 当 \(0 < \|x\| < \frac{1}{2}\) 时 \(g(x) = x\), 且 g 在 \(\|x\| \geqslant \frac{1}{2}\)\(\operatorname{Diff}^{\infty}(D)\) 中充分接近 e. 令 f 如 (e) 中, 对 \textless, 令 {}\textless{}\(f_{i}: T \to T\) 如下定义: 对于 \(x \in T\), \(y \in p^{-1}(x)\), \(z \in p^{-1}(f(x))\) 且 \(|z - y| < \pi\), 记 \(f_{i}(x) = p(ty + (1 - t)z)\). 则存在 \(h \in \operatorname{Diff}^{\infty}(Y)\), 使得对 \(h(x, y) = (f_{2\|x\|}(x), g(y))\)\(x \in T\), \(y \in D\) 有 , 且对  有 \(h(u) = u\); 此外, 可以选择 f 和 g, 使 h 在 \(u \in Y - (T \times D)\)\(\operatorname{Diff}^{\infty}(Y)\) 中任意接近 e. \(T \times \{0\}\) 中的点正是这样的 \(y \in Y\): 当 n 趋于  时, \(h^{kn}(x)\) 趋于一个不被 h 固定的点 (事实上它是周期为 k 的周期点). 因此, 每个与 h 交换的 Y 到自身的同胚都保持 \(T \times \{0\}\) 不变. 然后应用 (e).
\item
  设 Y 是维数 \(\geqslant1\) 的紧可微流形. 不存在李群 G 以及连续态射 \(\operatorname{Diff}^{\infty}(Y)\to G\), \(G\to\operatorname{Diff}^{0}(Y)\), 使它们的复合为 \(\operatorname{Diff}^{\infty}(Y)\) 到 \(\operatorname{Diff}^{0}(Y)\) 的典范嵌入. (使用 (d).)
\end{enumerate}

\begin{enumerate}
\def\labelenumi{\arabic{enumi}.}
\setcounter{enumi}{7}
\tightlist
\item
  设 G 是有限维李群. 假设 L(G) 是单的. 设 A 是 G 的正规子群. 若 A 不是开的, 则 A 是离散的, 且其在 G 中的中心化子是开的. (令 a 为 A 在 e 处相切的子代数. 证明 \(a = \{0\}\). 取 A 中不属于 G 中心的元素 x. 考虑从 G 到 A 的映射 \(y \mapsto xyj^{-1}\). 由 \(a = \{0\}\) 推出 G 中 x 的中心化子是开的. 然后利用指数映射证明, x 在 G 中的中心化子包含一个与 x 无关的 e 邻域.)
\end{enumerate}

¶ 9. 设 G 是 K 上维数 n 的紧李群.~假设李代数 L(G) 是单的. 对所有 \(g \in G\), 每个整数 \(m \geqslant 1\) 以及 G 中 e 的每个邻域 V, 令 M(g, m, V) 表示 G 中形如 \(\prod_{i=1} x_i(y_i, g) x_i^{-1}\) 的元素的集合, 其中 \(x_1, \ldots, x_m, y_1, \ldots, y_m\) 属于 V.

\begin{enumerate}
\def\labelenumi{(\alph{enumi})}
\tightlist
\item
  证明若 g 是中心化子不开的 G 中元素, m 是一个 \(g\)\(G\)\(m\)\(\geqslant n\) 的整数且 V 是 e 的邻域, 则 M(g, m, V) 是 e 的邻域. (存在 \(V\)\(e\)\(M(g, m, V)\)\(e\)\(a \in L(G)\), 使 \((\operatorname{Ad} g)(a) \neq a\). 令 \(\phi\) 为 G 的一个指数映射. 令 y 为 0 处映射 \(G\)\(y\)\(\lambda \mapsto \langle \phi(xa), g \rangle\) 的切映射对 1 的像 (其中 \(\lambda \in K\)). 则 \(b = (\operatorname{Ad} g - 1)(a) \neq 0\). 存在 V 中的 \(x_1, \ldots, x_m\), 使 \(V\) 包含 L(G) 的一个基. 考虑映射\(\{(\operatorname{Ad} x_1)(b), \ldots, (\operatorname{Ad} x_m)(b)\}\)\(L(G)\)
\end{enumerate}

\[
\left(x _ {1} ^ {\prime}, \dots , x _ {m} ^ {\prime}, y _ {1}, \dots , y _ {m}\right) \mapsto \prod_ {i = 1} ^ {n} x _ {i} ^ {\prime} (y _ {i}, g) x _ {i} ^ {\prime - 1}
\]

从 \(\mathbf{G}^{2m}\) 到 G. 证明在 \((x_{1},\ldots ,x_{m},e,\ldots ,e)\) 处的切映射是满射.)

\begin{enumerate}
\def\labelenumi{(\alph{enumi})}
\setcounter{enumi}{1}
\item
  令 \(\mathbf{U}\) 为 \(e\)\(\mathbf{G}\) 中 e 的一个邻域. 令 m 为一个 \(m\)\(\geqslant 1\) 的整数. 证明存在一个中心化子不开的 \(g\)\(\mathbf{G}\) 中元素 g, 使得 \(\mathbf{M}(g, m, \mathbf{G}) \subset \mathbf{U}\). (利用 \(\mathbf{G}\) 的紧性作反证.)
\item
  设 \(\mathbf{G}'\) 是其李代数为单的 K 上的紧李群. 若 \(\mathbf{K}\)\(\phi: \mathbf{G} \to \mathbf{G}'\) 是抽象群的双射同态, 则 \(\phi\) 连续. (对 \(\mathbf{G}'\) 应用 (b), 对 \(\mathbf{G}\) 应用 (a).)
\end{enumerate}

\begin{enumerate}
\def\labelenumi{\arabic{enumi}.}
\setcounter{enumi}{9}
\item
  设 G 是 K 上的李群. 证明存在 e 的一个邻域 V, 具有如下性质: 对 V 中元素的任意序列 \(G\)\(K\)\(V\)\(e\)\((x_{n})\), 若归纳定义 \(V\)\(y_{1} = x_{1}, y_{n} = (x_{n}, y_{n-1})\), 则序列 \((y_{n})\) 趋于 e.\(e\)
\item
  \begin{enumerate}
  \def\labelenumii{(\alph{enumii})}
  \tightlist
  \item
    对于  中的 \(x, y\), 定义 \(\mathbf{S}_{2n-1} \subset \mathbf{C}^n\)\(\alpha(x, y) \in [0, \pi]\) 为
  \end{enumerate}
\end{enumerate}

\[
\cos \alpha (x, y) = \mathscr {R} ((x | y)).
\]

对于  中的 \(s, t\), 记 . 证明 d 是 \(\mathbf{U}(n)\) 上左右不变的距离.\(d(s, t) = \sup_{x \in \mathbf{S}_{2n-1}} \alpha(sx, tx)\)\(d\)\(\mathbf{U}(n)\)

\begin{enumerate}
\def\labelenumi{(\alph{enumi})}
\setcounter{enumi}{1}
\item
  设 \(s \in \mathbf{U}(n)\). 令 \(\theta_1, \ldots, \theta_m\) 为 \(\mathbf{J} - \pi, \pi \mathbf{j}\) 中的数, 使 \(e^{i\theta_j}\) 是 s 的互异特征值. 令 \(s\)\(\theta(s) = \sup_{1 \leqslant j \leqslant m} |\theta_j|\). 则 \(d(e, s) = \theta(s)\). (令 \(V_j\) 为 s 对应于 \(s\)\(e^{i\theta_j}\) 的特征子空间. 对 \(x \in S_{2n-1}\), 通过按  分解, 求 \(\mathcal{R}((x|sx))\) 的一个下界.)\(V_j\)
\item
  设 \(s, t\), 满足 \(\mathbf{U}(n)\)\(\theta(t) < \frac{\pi}{2}\), 且 s 与 \(s\)\((s, t)\) 交换. 则 s 与 t 交换. (在 (b) 的记号下, 令 \(s\)\(t\)\(\mathrm{V}_j'\) 为 \(\mathrm{V}_j\) 的正交补, \(\mathrm{W}_j = t(\mathrm{V}_j)\); 证明 \(\mathrm{W}_j = (\mathrm{W}_j \cap \mathrm{V}_j) + (\mathrm{W}_j \cap \mathrm{V}_j')\), 然后证明 \(\mathrm{W}_j \cap \mathrm{V}_j' = \{0\}\).)
\item
  证明存在  中 e 的一个紧邻域 \(\mathbf{V}\), 它具有练习 10 的性质, 在 \(e\)\(\mathbf{U}(n)\) 的内自同构下不变, 并且若 x, y 是 \(\mathbf{U}(n)\)\(x, y\)\(\mathbf{V}\) 中不可交换的元素, 则 x 与 \(x\)\((x, y)\) 不可交换. (使用 (c).)
\item
  令 \(\beta\) 为 \(\mathbf{U}(n)\) 的归一化 Haar 测度. 令 \(\mathrm{V}_1\) 是 \(e\)\(\mathbf{U}(n)\) 中 e 的对称紧邻域, 使 \(\mathrm{V}_1^2 \in \mathrm{V}\). 令 p 为满足 \(p\) 的整数. 证明对每个 \(\beta (\mathrm{V}_1) > 1 / p\)\(\mathbf{GL}(n,\mathbf{C})\) 的有限子群 F, 都存在 F 的交换正规子群 A, 使 \(\operatorname{Card}(\mathrm{F} / \mathrm{A}) \leqslant p\) (Jordan 定理). (将其归约为 \(\mathrm{F} \subset \mathbf{U}(n)\) 的情形. 取 A 为由 \(\mathrm{F} \cap \mathrm{V}\) 生成的 F 的子群, 并使用 (d).)
\end{enumerate}

\begin{enumerate}
\def\labelenumi{\arabic{enumi}.}
\setcounter{enumi}{11}
\tightlist
\item
  设 G 是有限维连通实或复李群, \(\alpha\) 是 G 上 p 次左不变微分形式. \(\alpha\) 同时右不变, 当且仅当 \(d\alpha = 0\). (注意右不变性的条件可写成
\end{enumerate}

\[
\wedge^ {p} (^ {t} \mathrm{Ad} (s)). \alpha (e) = \alpha (e)
\]

对所有 \(s \in G\) 都成立, 因而等价于条件

\[
\sum_ {j = 1} ^ {f} \left\langle \alpha (e), u _ {1} \wedge \dots \wedge u _ {j - 1} \wedge [ u, u _ {j} ] \wedge u _ {j + 1} \wedge \dots \wedge u _ {p} \right\rangle = 0
\]

对 \(u, u_1, \ldots, u_p\)\(\mathrm{L}(\mathrm{G})\) 中的所有  都成立. 特别地, \(\mathbf{G}\) 上 1 次左右不变微分形式构成维数为

\[
\dim \mathrm{L} (\mathrm{G}) - \dim [ \mathrm{L} (\mathrm{G}), \mathrm{L} (\mathrm{G}) ].
\]

\begin{enumerate}
\def\labelenumi{\arabic{enumi}.}
\setcounter{enumi}{12}
\tightlist
\item
  令 \(\mathbf{R}^n\) 带有通常的内积 \(((\xi_i), (\eta_i)) \mapsto \sum_{i=1} \xi_i \eta_i\). 令 \(\mathbf{I}(n)\) 为 \(\mathbf{R}^n\) 的等距仿射变换群. 它典范地同构于 \(\mathbf{GL}(n + 1, \mathbf{R})\) 的李子群, 该子群由矩阵
\end{enumerate}

\[
S = \left( \begin{array}{c c} U & x \\ 0 & 1 \end{array} \right)
\]

其中 \(U \in \mathbf{O}(n)\), x 是 \(x\)\(\mathbf{R}^n\) 的任意元素 (类型为 \((n,1)\) 的矩阵). \(\mathbf{I}(n)\) 可识别为由 \(\mathbf{O}(n)\) 到  的典范嵌入所定义的  与 \(\mathbf{R}^n\) 的半直积 (此时矩阵 S 记作 \(\mathbf{O}(n)\)\(\mathbf{GL}(n,\mathbf{R})\)\(S\)\((U,x)\)). 令 \(p\colon \mathbf{I}(n)\to \mathbf{O}(n)\) 为典范满射, 核为 \(\mathbf{R}^n\).

\begin{enumerate}
\def\labelenumi{(\alph{enumi})}
\tightlist
\item
  设 \(\Gamma\) 是 \(\mathrm{I}(n)\) 的离散子群. 考虑 \(\tilde{p}(\Gamma)\)\(p(\Gamma)\) 在 \(\mathbf{O}(n)\) 中的闭包 \(V\). 证明其单位元连通分支是交换的. (令 V 为 \(I\)\(\mathbf{O}(n)\) 中 I 的一个邻域, 具有练习 11 (d) 的性质, 并且还满足对所有 \(d\) 有 \(\|U - I\| \leqslant \frac{1}{4}\) (\(U \in V\)\(\operatorname{End}(\mathbf{R}^n)\) 上的范数由 \(\mathbf{R}^n\) 上的欧氏范数导出). 反证, 假设存在 \(S_1 = (U_1, x_1)\), \(S_2 = (U_2, x_2)\), 其中 \(U_1\) 和 \(U_2\) 属于 V 且不交换. 证明 \(V\)\((S_1, S_2) = ((U_1, U_2), y)\), 其中
\end{enumerate}

\[
\| y \| \leqslant \frac {1}{4} (\| x _ {1} \| + \| x _ {2} \|).
\]

\(S_{k}\). 然后对  归纳定义 \(S_{k}=(S_{1},S_{k-1})\). 利用 V 的选取, 证明此序列有无穷多个不同的项, 且在 \(k\geqslant3\) 中有界, 矛盾.)\(\mathbf{I}(n)\)

\begin{enumerate}
\def\labelenumi{(\alph{enumi})}
\setcounter{enumi}{1}
\item
  若 \(\Gamma\) 是晶体学群, 即 \(\mathrm{I}(n)/\Gamma\) 是紧的, 则证明对所有 \(x \in \mathbf{R}^n\), 由  生成的  的仿射子空间 \(\mathbf{L}\) 等于 \(\mathbf{R}^n\)\(\Gamma x\)\(\mathbf{R}^n\). (否则, 对所有 \(y \in \mathbf{R}^n\), \(\Gamma y\) 的所有点到 \(\mathbf{L}\) 的距离都相同; 由于这个距离可以任意大, \(\mathrm{I}(n)/\Gamma\) 不是紧的.)
\item
  若 \(\Gamma\) 是交换的晶体学群, 则 \(\Gamma \subset \mathbf{R}^n\). (若 \(S = (U, x) \in \Gamma\), 使得 \(\mathrm{V} \subset \mathbf{R}^n\) 中由 U 的不动点组成的向量子空间 \(U\) 不等于 \(\mathbf{R}^n\), 且 \(\mathrm{V}^\perp\) 表示 \(\mathrm{V}\) 的正交子空间, 则对所有 \(S' = (U', x') \in \Gamma\), U' 都保持 \(U'\)\(\mathrm{V}\) 和 \(\mathrm{V}^\perp\) 不变. 另一方面, 由于 \(U|\mathrm{V}^\perp\) 没有非零不动点, 证明改变原点后可设 \(x \in \mathrm{V}\). 根据 (b), 存在 \(S' = (U', x') \in \Gamma\), 使 x' 在 \(y'\)\(x'\)\(\mathrm{V}^\perp\) 上的正交投影 \(\neq 0\). 用公式 \(Sx'\)\(S = S'SS'^{-1}\) 计算 Sx', 推出 \(U.y' = y'\), 这与 \(\mathrm{V}\) 的定义矛盾.
\item
  若 \(\Gamma\) 是晶体学群, 则 \(\Gamma \cap \mathbf{R}^n\) 是秩为 n 的自由交换群, 且 \(n\)\(\Gamma / (\Gamma \cap \mathbf{R}^m)\) 是有限群 (Bieberbach 定理). (令 W 为由 \(W\) 生成的 \(\mathbf{R}^n\) 的向量子空间. 紧群 \(\Gamma \cap \mathbf{R}^n\)\(p(\Gamma)\) 只有有限个连通分支. 若 \(W = \{0\}\), 则由 (a), \(\Gamma\) 含有有限指数的交换正规子群 \((a)\)\(\Gamma_1\); 此时 \(\Gamma_1\) 也是晶体学群, 这与 (c) 矛盾. 因而 \((c)\)\(W \neq \{0\}\). 证明 \(p(\Gamma)\) 保持 W 不变且 \(W\)\(p(\Gamma)|W\) 是有限的; 否则存在  中的序列 \((S_m)\), 若 \(\Gamma\)\(S_m = (U_m, x_m)\), 则 \(U_m\) 彼此不同且趋于 I; 然后构造 \(I\)\((I, a_j)S_m(I, a_j)^{-1s_m^1}\), 其中 \((a_j)\) 是 \(\Gamma \cap \mathbf{R}^n\) 的一个基, 这将与 \(\Gamma\) 离散的假设矛盾. 最后, 为了看出 \(W = \mathbf{R}^n\), 证明否则 \(\Gamma\) 在 \(\mathbf{R}^n/W\) 上的作用将是一个不含非零平移的晶体学群的作用.)\(\neq 0\)
\end{enumerate}

\exercisesfor{5}

\begin{enumerate}
\def\labelenumi{\arabic{enumi}.}
\tightlist
\item
  假设 \(\mathbf{K} = \mathbf{R}\). 设 \(r\) 是一个 \(\geqslant 1\) 的整数或 \(\infty\). C＾r 类的群是一个集合 \(\mathbf{C}^r\)\(\mathbf{G}\), 它具有群结构和 \(\mathbf{C}^r\) 类流形结构, 并且从 \((x,y)\mapsto xy^{-1}\)\(\mathbf{G}\times \mathbf{G}\) 到 \(\mathbf{G}\) 的映射  是 \(\mathbf{C}^r\) 类的.
\end{enumerate}

\begin{enumerate}
\def\labelenumi{(\alph{enumi})}
\tightlist
\item
  以下设 \(r \geqslant 2\). 将 e 的一个邻域识别为 Banach 空间的一个邻域. 令 \(e\)\(G\)\(xy = P(x, y)\), 其中 P 是在 \(P\)\(C^2\)\((0, 0)\) 的一个开邻域上定义的 C＾2 类映射. 令
\end{enumerate}

\[
\left(\mathrm{D} _ {1} \mathrm{D} _ {2} \mathrm{P}\right) (0, 0) = \mathrm{B}.
\]

证明

\[
x y = x + y + \mathrm{B} (x, y) + | x | | y | o (1)
\]

当 \((x,y)\to (0,0)\) 时成立. (使用带积分余项的二阶展开.)

\begin{enumerate}
\def\labelenumi{(\alph{enumi})}
\setcounter{enumi}{1}
\tightlist
\item
证明
\end{enumerate}

\[
\begin{array}{c} x ^ {- 1} = - x + \mathrm{B} (x, x) + | x | ^ {2} o (1) \\ x y x ^ {- 1} = y + \mathrm{B} (x, y) - \mathrm{B} (y, x) + | x | | y | o (1) \\ x ^ {- 1} y ^ {- 1} x y = \mathrm{B} (x, y) - \mathrm{B} (y, x) + | x | | y | o (1) \end{array}
\]

当 \((x,y)\) 趋于 \((0,0)\) 时成立.

\begin{enumerate}
\def\labelenumi{\arabic{enumi}.}
\setcounter{enumi}{1}
\tightlist
\item
  设 G 是一个 C＾r 类群, 其中 \(G\)\(C^r\)\(r \geqslant 2\).
\end{enumerate}

\begin{enumerate}
\def\labelenumi{(\alph{enumi})}
\item
  设 \(t \in \mathcal{T}^{(s)}(\mathrm{G})\), \(t' \in \mathcal{T}^{(s')}(\mathrm{G})\). 若 \(s + s' \leqslant r\), 按照 § 3 第 1 号定义 \(t * t' \in \mathcal{T}^{(s + s')}(\mathrm{G})\). 若 t 和 t' 都没有常数项, 则 t * t' 也没有常数项. t 在从 \(t\)\(t'\)\(t * t'\)\(t\)\(x \mapsto x^{-1}\)\(\mathrm{G}\) 到 \(\mathrm{G}\) 的映射  下的像记为 \(t^{\vee}\); 于是 \(t^{\vee} \in \mathcal{T}^{(s)}(\mathrm{G})\). 若 f 是 \(f\) 上取值于 Hausdorff 多赋范空间的 \(\mathrm{C}^r\) 类函数, 则 f * t 表示 \(\mathrm{G}\)\(f * t\)\(\mathrm{G}\) 上定义为 \((f * t)(x) = \langle \varepsilon_x * t'^{\nu}, f \rangle\) 的函数, 对所有 \(x \in \mathrm{G}\) 都如此; 当 \(\mathrm{C}^r - s\)\(s < \infty\) 时, 该函数是  类的. 证明仿照 § 3 第 4 号可得 \(f * (t * t') = (f * t) * t'\).
\item
  若 \(t \in \mathrm{T}_e(\mathrm{G})\), 则  上的向量场 \(x \mapsto \varepsilon_x * t\) 记为 \(\mathbf{G}\)\(\mathbf{L}_t\). 证明仿照 § 3 第 6 号, 对于 \(t, t'\)\(\mathrm{T}_e(\mathrm{G})\) 中的 t, t', 有 \(\mathrm{L}_{t^* t'} = \mathrm{L}_t \circ \mathrm{L}_{t'}\), 因而
\end{enumerate}

\[
\mathrm{L} _ {t ^ {*} t ^ {\prime} - t ^ {*} t} = \left[ \mathrm{L} _ {t}, \mathrm{L} _ {t ^ {\prime}} \right].
\]

\[
t * t ^ {\prime} - t ^ {\prime} * t
\]

\[
\mathrm{T} _ {e} (\mathrm{G})
\]

\begin{enumerate}
\def\labelenumi{(\alph{enumi})}
\setcounter{enumi}{2}
\tightlist
\item
  使用练习 1 (a) 的记号, 若 \(t \in \mathrm{T}_e(\mathrm{G})\), 则
\end{enumerate}

\[
[ t, t ^ {\prime} ]
\]

\[
\mathrm{L} _ {t} (x) = (\mathrm{D} _ {2} \mathrm{P} (x, 0)) (t)
\]

当 x 充分接近 0 时成立. 推出 \(x\)\([t, t'] = \mathbf{B}(t, t') - \mathbf{B}(t', t)\).

\begin{enumerate}
\def\labelenumi{(\alph{enumi})}
\setcounter{enumi}{3}
\tightlist
\item
  证明 \(\mathrm{T}_e(\mathrm{G})\), 取括号 \((t, t') \mapsto [t, t']\), 是一个可赋范李代数. (证明 Jacobi 恒等式时, 使用 (c)、练习 1 (b) 以及《代数》第 I 章 §6 第 2 号的恒等式 (5).)
\end{enumerate}

(关于本练习的后续内容, 参见 § 8, 练习 6.)

\exercisesfor{6}

\begin{enumerate}
\def\labelenumi{\arabic{enumi}.}
\item
  令 D 为 \(\mathbf{SL}(2,\mathbf{C})\) 中形如 \(\begin{pmatrix}a & b \\ 0 & a^{-1}\end{pmatrix}\) 的元素集合, 其中 a \textgreater{} 0 且 \(b \in C\). 证明 D 是 \(\mathbf{SL}(2,\mathbf{C})\) 的底层实李群的李子群, 并且映射 \((u,d) \mapsto ud\) 从 \(\mathbf{SU}(2,\mathbf{C}) \times \mathbf{D}\) 到 \(\mathbf{SL}(2,\mathbf{C})\) 是实解析流形的同构. 由此以及 §3 练习 7 (c) 推出 \(\mathbf{SL}(2,\mathbf{C})\) 是单连通的.
\item
  令 G 为 \(G\)\(\mathbf{SL}(2, \mathbf{R})\) 的万有覆盖, \(\pi\) 为从 G 到 \(G\)\(\mathbf{SL}(2, \mathbf{R})\) 的典范态射, N 为其核.\(N\)
\end{enumerate}

\begin{enumerate}
\def\labelenumi{(\alph{enumi})}
\item
  仿照练习 1, 证明存在从实解析流形 \(\mathbf{U} \times \mathbf{R}^2\) 到实解析流形 \(\mathbf{SL}(2, \mathbf{R})\) 的同构. 推出 N 同构于 \(N\)\(\mathbf{Z}\).
\item
  若 \(\rho\) 是 G 在复向量空间上的解析线性表示, 则 \(G\)\(\operatorname{Ker} \rho \supset N\). ( 的表示 \(\mathrm{L}(\rho)\) 通过复化定义出 \(\mathfrak{sl}(2, \mathbf{R})\)\(\mathfrak{sl}(2, \mathbf{C})\) 的一个表示, 后者根据练习 1 具有 \(\mathrm{L}(\sigma)\) 的形式, 其中 \(\sigma\) 是 \(\mathbf{SL}(2, \mathbf{C})\) 的一个解析线性表示. 然后 \(\mathrm{L}(\sigma \circ \pi) = \mathrm{L}(\rho)\), 因而 \(\sigma \circ \pi = \rho\).)
\end{enumerate}

\begin{enumerate}
\def\labelenumi{\arabic{enumi}.}
\setcounter{enumi}{2}
\item
  令 \(G_{1}\) 为 § 4 练习 5 中定义的单连通幂零实李群. 令 Z 为 \(G_{1}\) 的中心, N 为 Z 的非平凡离散子群, 且 \(G = G_{1}/N\). 若 \(\rho\) 是 G 的有限维解析线性表示, 则 \(\rho(Z/N)\) 是一组半单自同构, 因为 Z/N 是紧的; 但 \(Z = (G, G)\), 因而 \(\rho(Z/N)\) 是幂零的 (第 I 章 § 6 命题 6). 所以 \(\rho\) 在 Z/N 上是平凡的.
\item
  设 G 是紧连通复李群. 证明 G 的每个解析线性表示都是平凡的.
\item
  在 § 1 练习 9 的记号下, 证明 \(\mathbf{H}'\) 是 H 的积分子群. 推出在单连通群 \(\mathbf{SU}(2,\mathbf{C})\times \mathbf{SU}(2,\mathbf{C})\) (§ 3, 练习 7 (c)) 中存在不是闭的单参数子群.
\end{enumerate}

¶ 6. 令 I 为带离散拓扑的区间 \(\{1,2\}\). 令 E 为函数 \(f: I \to R\) 所成的完备赋范空间, 其中 \(\|f\| = \sum_{i \in I} |f(i)| < +\infty\). 对所有 \(x \in I\), 令 \(\varepsilon_x\) 为 E 中满足当  时 \(\varepsilon_x(t) = 0\) 且 \(t \neq x\)\(\varepsilon_x(x) = 1\) 的元素. 令 P 为由 \(x\varepsilon_x\) (\(x \in I\)) 生成的 E 的子群; 它是 E 的离散子群. 令 G 为李群 E/P. 令 F 为 E 中由满足  的 \(f \in E\) 组成的超平面. 令 H 为李群 F/(F \(\sum_{i \in I} f(i) = 0\) P). 令 \(\phi\) 为从 H 到 G 的典范态射. 证明 \(\phi\) 是双射且是浸入, 但不是李群同构. (令 f 为 E 上核为 F 的线性形式; 证明 \(f(P) = R\), 从而 \(F + P = E\).) 推出命题 3 中的可数性假设不可省略.

\begin{enumerate}
\def\labelenumi{\arabic{enumi}.}
\setcounter{enumi}{6}
\tightlist
\item
  令 \(\phi\) 为从 \(\mathbf{Z} / 2\mathbf{Z}\) 到 \(\mathbf{C}\) 的置换群的态射, 将非单位元映到 \(z\mapsto \overline{z}\). 令 \(\mathrm{G}\) 为由  相应的 \(\mathbf{Z} / 2\mathbf{Z}\) 与实李群 \(\mathbf{C}\) 构成的半直积. 它是一个李代数为 \(\phi\)\(\mathbf{R}^2\) 的实李群. 证明 \(\mathrm{G}\) 上不存在与其实李群结构相容的复李群结构.
\end{enumerate}

¶ 8. 设 H 是维数 \(\aleph_{0}\) 的复 Hilbert 空间, G 是视为实李群的 H 的酉群. 群 G 是单连通的. \(\dagger\) (a) 令 Z 为 G 的中心, 它同构于 T. 令 a 为无理数. 令 \(\hat{\mathbf{a}}\) 为 \(\mathrm{L}(\mathrm{G})\times \mathrm{L}(\mathrm{G})\) 中由 \((x,ax)\) (\(x\in \mathrm{L}(Z)\)) 组成的李子代数. 令 S 为相应的 \(\mathrm{G}\times \mathrm{G}\) 的积分子群. 则 \(\hat{\mathbf{a}}\) 是 \(\mathrm{L}(\mathrm{G})\times \mathrm{L}(\mathrm{G})\) 的理想; 然而 S 不是闭的, 且在 \(\mathbf{Z}\times \mathbf{Z}\) 中稠密.

\begin{enumerate}
\def\labelenumi{(\alph{enumi})}
\setcounter{enumi}{1}
\tightlist
\item
  令 \(\mathfrak{g} = (\mathrm{L}(\mathrm{G})\times \mathrm{L}(\mathrm{G})) / \tilde{\mathfrak{s}}\). 不存在李代数为 \(\mathfrak{g}\) 的李群. (若 \(\mathrm{H}\) 是这样的群. 由于 \(\mathrm{G}\) 是单连通的, 存在态射 \(\phi : \mathrm{G}\times \mathrm{G}\to \mathrm{H}\), 使 \(\mathrm{L}(\phi)\) 是从 \(\mathrm{L}(\mathrm{G})\times \mathrm{L}(\mathrm{G})\) 到 \(\mathfrak{g}\) 的典范映射. 令 \(\mathrm{N} = \operatorname{Ker}\phi\). 则 \(\mathrm{L}(\mathrm{N})\supset \tilde{\mathfrak{s}}\), 因而 \(\mathrm{N}\supset \mathrm{S}\), 再由 (a) 得 \(\mathrm{N}\supset \mathrm{Z}\times \mathrm{Z}\). 于是 \(\mathrm{L}(\mathrm{N})\supset \mathrm{L}(\mathrm{Z})\times \mathrm{L}(\mathrm{Z})\), 矛盾.)
\end{enumerate}

\begin{enumerate}
\def\labelenumi{\arabic{enumi}.}
\setcounter{enumi}{8}
\item
  设 X 是维数为 n 的非空连通紧复流形.~假设 X 上存在全纯向量场 \(\xi_{1},\ldots,\xi_{n}\), 它们在 X 的每一点都线性无关. 证明存在复李群 G 及其离散子群 D, 使 X 微分同胚于 G/D. (令 \(c_{ijk}\) 为 X 上满足 \([\xi_{i},\xi_{j}]=\sum_{k}c_{ijk}\xi_{k}\) 的全纯函数. 因为 X 紧, \(c_{ijk}\) 是常数. 取 G 为李代数以 \(c_{ijk}\) 为结构常数的单连通复李群, 并使用定理 5.)
\item
  设 G 是具有有限个连通分支的李群. G 是幺模的, 当且仅当对所有  都有 Tr ad \(a = 0\).\(a \in \mathrm{L}(\mathrm{G})\)
\item
  设 E 是 \(\mathbf{C}\) 上的完备可赋范空间, \(v \in \mathcal{L}(\mathrm{E})\), 且 \(g = \exp(v)\). 假设 \(\operatorname{Sp}(v) \cap 2i\pi(\mathbf{Z} - \{0\}) = \emptyset\). 令 \(\mathrm{E}_1, \mathrm{E}_2\) 为 E 的在 v 下不变的闭向量子空间, 且 \(v\)\(\mathrm{E}_2 \subset \mathrm{E}_1\). 假设由 g 在 \(\mathrm{E}_1 / \mathrm{E}_2\) 上导出的自同构是恒等映射. 则 \(g\)\(v(\mathrm{E}_1) \subset \mathrm{E}_2\).
\end{enumerate}

¶ 12. 考虑实或复完备可赋范空间 E, 以及 E 的一个闭向量子空间 F, 满足 \(F^{0}\) 在 \(E^{\prime}\dagger\) 中没有拓扑补空间 (a) 令 A (分别为 B) 为 E (分别为 F) 的连续自同态所成的完备可赋范空间. 令 C 为满足  的 \(u \in A\) 的集合. 令 \(u(F) \subset F\)\(\alpha\) 为从 C 到  的映射 \(u \mapsto (u, u \mid F)\). 则 \(A \times B\)\(\alpha\) 是从完备可赋范空间 C 到 \(A \times B\) 的一个闭向量子空间上的同构, 且 \(\alpha(C)\) 在 \(A \times B\) 中没有拓扑补空间. (令 \(x \in E\) 满足 \(x \notin F\). 令 \(\xi \in \mathrm{F}^0\) 满足 \(\langle x, \xi \rangle = 1\). 若 \(\eta \in \mathrm{E}'\), 令 \(\tilde{\eta}\) 表示 A 中的元素 \(y \mapsto \langle y, \eta \rangle x\). 假设存在从 A × B 到  的投影算子 \(\pi\). 定义 \(\alpha(\mathrm{C})\)\(\tilde{\pi} \colon \mathrm{E}' \to \mathrm{E}'\) 为 \(\tilde{\pi}(\eta) = t(\alpha^{-1}(\pi(\tilde{\eta}, 0)))(\xi)\). 则 \(\tilde{\pi}\) 是从 \(\mathrm{E}'\) 到 \(\mathrm{F}^0\) 的投影算子, 矛盾.)

如下: \(M = \bigoplus_{i=0}^{n} M_{i}\) 按 \(M_{i}\) 分次; \(M_{0}\) 是 1 的标量倍数的集合; \(M_{i}\) 是非零元素 u 的标量倍数的集合; \(M_{2} = \{0\}\); \(M_{3} = F\); \(M_{4} = E\); 当 i \(M_{i} = \{0\}\)\textgreater{} 5 时 ; 对所有 \(x \in M_{3}\) 有 ux = x. 令 N 为完备可赋范空间 M 的连续自同态所成的完备可赋范空间. 令 \(N_{1}\) 为 M 的连续导子的集合. 利用 (a) 证明 \(N_{1}\) 在 N 中没有拓扑补空间.

\begin{enumerate}
\def\labelenumi{(\alph{enumi})}
\setcounter{enumi}{2}
\tightlist
\item
  证明 M 的双连续自同构群是 \(\mathbf{GL}(\mathbf{M})\) 的李拟子群, 但不是 \(\mathbf{GL}(\mathbf{M})\) 的李子群. (使用命题 18 的推论 2 以及 (b).)
\end{enumerate}

\begin{enumerate}
\def\labelenumi{\arabic{enumi}.}
\setcounter{enumi}{12}
\item
  设 \(\mathbf{G}\) 是实李群, \(\mathbf{L} = \mathbf{L}(\mathbf{G})\), 且 \(\phi: \mathbf{L} \to \mathbf{G}\) 是在 0 处可微的映射, 满足 \(\mathrm{T}_0(\phi) = \mathrm{Id}_{\mathrm{L}}\), 并且对所有 \(\phi(nx) = \phi(x)^n\)\(x \in \mathbf{L}\) 和 \(n \in \mathbf{Z}\) 有 . 则 \(\phi = \exp_0\). (令 \(\mathbf{V}\) 为 \(\mathbf{L}\) 中 0 的一个邻域, \(\mathbf{W}\) 为 \(e\)\(\mathbf{G}\) 中 e 的一个邻域, 使得 \(0 = \exp_0 | \mathbf{V}\) 是从 \(\mathbf{V}\) 到 \(\mathbf{W}\) 的解析同构. 令 \(\psi = \theta^{-1} \circ (\phi | \phi^{-1}(\mathbf{W}))\). 则 \(\mathrm{T}_0(\psi) = \mathrm{Id}_{\mathrm{L}}\), 从而利用当 x 充分接近 0 且  时的等式 , 得 \(\psi = \mathrm{Id}_{\mathrm{L}}\).\(\psi\left(\frac{1}{n}x\right) = \frac{1}{n}\psi(x)\)\(x\)\(n \in \mathbf{Z} - \{0\}\)
\item
  令 \(\mathfrak{a}_1\) 为交换李代数 \(\mathbf{R}^4\), 将其识别为 \(\mathbf{C}^2\). 令 \(\mathfrak{a}_2\) 为交换李代数 \(\mathbf{R}\). 令 \(\phi\) 为从 \(\mathfrak{a}_2\) 到 \(\mathrm{Der}(\mathfrak{a}_1)\) 的同态, 将 1 映到导子 \((z_1, z_2) \mapsto (iz_1, i\sqrt{2}z_2)\). 令 \(\mathfrak{a}\) 为由  相应的 \(\mathfrak{a}_2\) 与 \(\mathfrak{a}_1\) 的半直积.\(\phi\)
\end{enumerate}

\begin{enumerate}
\def\labelenumi{(\alph{enumi})}
\tightlist
\item
  证明 (Int a) \textbar{} \(a_{1}\) 对所有 \(\phi \in R\) 都包含自同构
\end{enumerate}

\[
\left(z _ {1}, z _ {2}\right) \mapsto \left(e ^ {i \phi} z _ {1}, e ^ {i \sqrt {2} \phi} z _ {2}\right).
\]

\begin{enumerate}
\def\labelenumi{(\alph{enumi})}
\setcounter{enumi}{1}
\item
  证明 (Int a) \textbar{} a₁ 在 GL(a₁) 中的闭包 G 对所有 (φ, φ') ∈ R² 都包含自同构 (z₁, z₂) ↦ (eⁱΦz₁, eⁱΦ'z₂).
\item
  证明 \(\mathbf{L}(\mathbf{G})\) 包含  的自同态 \(u: (z_1, z_2) \mapsto (z_1, 0)\), 且不存在 \(\mathfrak{a}_1\)\(x \in \mathfrak{a}\) 使 \(u = (\operatorname{ad} x) \mid \mathfrak{a}_1\).
\item
  推出 \(\operatorname{Int}(\mathfrak{a})\) 在 \(\operatorname{Aut}(\mathfrak{a})\) 中不是闭的.
\end{enumerate}

\begin{enumerate}
\def\labelenumi{\arabic{enumi}.}
\setcounter{enumi}{14}
\tightlist
\item
  证明 §3 练习 7 (a) 的有限维单连通实李群包含非单连通的李子群 (事实上, 它包含同构于 U 的李群).
\end{enumerate}

  ¶ 16. (a) 设 g 是复李代数. 令 \(\bar{g}\) 为使用 C 的自同构 \(\lambda \mapsto \bar{\lambda}\) 从 g 导出的复李代数. 令 \(g_0\) 为实李代数
由限制标量从 g 导出. 令 \(g' \supset g_0\) 为复李代数 \(g_0 \otimes_{R} C\). 对所有 \(x \in g\), 记

\[
f (x) = \frac {1}{2} (x \otimes 1 - (i x) \otimes i) \in \mathfrak {g} ^ {\prime}, \quad g (x) = \frac {1}{2} (x \otimes 1 + (i x) \otimes i) \in \mathfrak {g} ^ {\prime}.
\]

证明 \(f(\text{resp. } g)\) 是从 \(\mathfrak{g}\) (分别从 \(\overline{\mathfrak{g}}\)) 到  的理想 \(\mathfrak{m}\) (分别为 \(\mathfrak{n}\)) 的同构, 且理想 \(\mathfrak{g}'\)\(\mathfrak{m}, \mathfrak{n}\) 在 \(\mathfrak{g}'\) 中互补. 这定义了从 \(\mathfrak{g}'\) 到 \(\mathfrak{m}, \mathfrak{n}\) 的投影算子. 证明对所有 \(x \in \mathfrak{g}, f(x) = p(x), g(x) = q(x)\), .

\begin{enumerate}
\def\labelenumi{(\alph{enumi})}
\setcounter{enumi}{1}
\item
  假设 g 是有限维的. 令 G 为李代数为 g 的单连通复李群. 令 \(\overline{G}\) 为共轭复李群, \(G_{0}\) 为底层实李群. 则 \(\mathrm{L}(\overline{\mathbf{G}})=\overline{\mathfrak{g}}\) 且 \(\mathrm{L}(G_{0})=g_{0}\). 令 M (分别为 N) 为李代数为 m (分别为 n) 的单连通复李群. 于是 f 定义了从 G 到 M 的同构 \(\phi\), g 定义了从  到 N 的同构 \(\gamma\), 而李代数为 \(\overline{G}\) 的单连通复李群 \(G'\) 与 M × N 对应. 证明 \(g'\)\(G'\) 中李代数为 \(g_{0}\) 的实积分子群是闭的、单连通的, 并与 \(G_{0}\) 对应. 将 \(pr_{1}\) (分别为 \(pr_{2}\)) 限制在 G  M × N 上, 得到从 G (分别从 ) 到 M (分别到 N) 的同构 \(\phi\) (分别为 \(\gamma\)).\(\overline{G}\)
\item
  由 (b) 以及第 10 号的注记 2 推出,  连同  到 \(\mathbf{G}'\) 的典范嵌入, 是 \(\mathbf{G}_0\) 的万有复化.\(\mathbf{G}'\)\(\mathbf{G}_0\)
\item
  设 H 是李代数为 g 的连通复李群, \(\bar{H}\) 是共轭复李群, \(H_{0}\) 是底层实李群, 且 G 是 H 的万有覆盖空间. 令 K 为 G 到 H 的典范态射的 (离散) 核. 证明 K 是 \(G'\) 的中心子群. 因而 \(G_{0}\) 到 \(G'\) 的典范嵌入定义了 \(H_{0}\) 到 \(G'/N\) 的嵌入 i. 证明 \((G'/N, i)\) 是 \(H_{0}\) 的万有复化. (注意这个有序对具有命题 20 的万有性质.)
\item
  记 \(G^{\prime}/N = (H_{0})_{\mathbf{c}}\). 由 (d) 推出从 \(\psi\)\((\mathrm{H}_{0})_{\mathbf{c}}\) 到 \(H \times \bar{H}\) 的典范态射 , 它是一个覆盖映射. 通过例子 \(H = C^{*}\) 证明 \(\psi\) 一般不是同构.
\end{enumerate}

\begin{enumerate}
\def\labelenumi{\arabic{enumi}.}
\setcounter{enumi}{16}
\tightlist
\item
  \begin{enumerate}
  \def\labelenumii{(\alph{enumii})}
  \tightlist
  \item
    令 G, \(\overline{G}\), \(G_{0}\) 如练习 16 (b) 中. 令 V 为有限维复向量空间, \(\rho\) 为 \(G_{0}\) 在 V 上的不可约线性表示. 证明存在有限维复向量空间 X, Y, G (分别为 ) 在 X (分别在 Y) 上的解析不可约线性表示 \(\sigma\) (分别为 \(\tau\)), 以及从 V 到 X \(\overline{G}\)\(\otimes\) Y 的同构, 将 \(\rho\) 变为 \(\sigma \otimes \tau\). (使用练习 16 (a)、第 I 章 § 2 的命题 2 以及《代数》第 VIII 章 § 7 的命题 8 的推论.)
  \end{enumerate}
\end{enumerate}

\begin{enumerate}
\def\labelenumi{(\alph{enumi})}
\setcounter{enumi}{1}
\tightlist
\item
  证明若去掉 G 单连通的假设, (a) 的结论不一定成立. (考虑复李群 \(\mathbf{C}^*\).)
\end{enumerate}

\begin{enumerate}
\def\labelenumi{\arabic{enumi}.}
\setcounter{enumi}{17}
\tightlist
\item
  设 \(\Lambda\) 是李代数为 \(a\) 的有限维连通交换复李群; 令 \(\Lambda\) 为 \(\exp_{A}\) 的核, 从而 \(\Lambda\) 与 \(a / \Lambda\) 对应.
\end{enumerate}

\begin{enumerate}
\def\labelenumi{(\alph{enumi})}
\tightlist
\item
  以下条件等价:
\end{enumerate}

(al) 典范映射 \(\mathbf{C} \otimes_{\mathbf{Z}} \Lambda \to a\) 是单射.

(a2) A 同构于某个 \((\mathbf{C}^{*})^{n}\) 的李子群.

(a3) A 同构于群 \((\mathbf{C}^{*})^{p}\times \mathbf{C}^{q}\).

(a4) A 具有有限维忠实复线性表示.

(a5) A 具有有限维半单忠实复线性表示, 且其像是闭的.

\begin{enumerate}
\def\labelenumi{(\alph{enumi})}
\setcounter{enumi}{1}
\tightlist
\item
  以下条件等价:
\end{enumerate}

(b1) 典范映射 \(\mathbf{C} \otimes_{\mathbf{Z}} \Lambda \to \mathfrak{a}\) 是满射.

(b2) A 同构于群 \((\mathbf{C}^{*})^{n}\) 的一个商.

(b3) A 的任何直因子都不同构于 C.

(b4) A 的每个复线性表示都是半单的.

\begin{enumerate}
\def\labelenumi{(\alph{enumi})}
\setcounter{enumi}{2}
\tightlist
\item
  以下条件等价:
\end{enumerate}

\begin{enumerate}
\def\labelenumi{(\roman{enumi})}
\setcounter{enumi}{149}
\tightlist
\item
  典范映射 \(\mathbf{C} \otimes_{\mathbf{Z}} \Lambda \to \mathfrak{a}\) 是双射.
\end{enumerate}

(c2) A 同构于群 \((\mathbf{C}^{*})^{n}\).

\begin{enumerate}
\def\labelenumi{(\alph{enumi})}
\setcounter{enumi}{3}
\tightlist
\item
  令 \(\mathbf{F}\) 为 \(\mathbf{A}\) 的有限子群, 令 \(\mathbf{A}' = \mathbf{A} / \mathbf{F}\). 证明 \(\mathbf{A}\) 满足条件 \((a_i)\) (分别为 \((b_i), (c_i)\)) 当且仅当 \(\mathbf{A}'\) 满足相应条件.
\end{enumerate}

\begin{enumerate}
\def\labelenumi{\arabic{enumi}.}
\setcounter{enumi}{18}
\tightlist
\item
  设 G 是实李群. 证明存在 L(G) 中 0 的一个邻域 V, 使得对 V 中的 \(G\)\(V\)\(L(G)\)\(x, y\) 有:\(V\)
\end{enumerate}

\[
\exp (t (x + y)) = \lim _ {n \rightarrow + \infty} \left(\exp \frac {t x}{n}. \exp \frac {t y}{n}\right) ^ {n}
\]

\[
\exp (t ^ {2} [ x, y ]) = \lim _ {n \rightarrow + \infty} \left(\exp \frac {t x}{n}. \exp \frac {t y}{n}. \exp \frac {- t x}{n}. \exp \frac {- t y}{n}\right) ^ {n ^ {2}}
\]

对 \(t \in (0, 1)\) 一致成立.

\begin{enumerate}
\def\labelenumi{\arabic{enumi}.}
\setcounter{enumi}{19}
\item
  设 G 是李群, H 是 G 的李子群, A 是 G 的积分子群, 且 \(\mathrm{L}(\mathrm{H}) \cap \mathrm{L}(\mathrm{A}) = \{0\}\). 证明 \(H \cap A\) 在李群 A 中是离散的.
\item
  设 \(G\) 是实李群, H 是正规的一维积分子群.\(H\)
\end{enumerate}

\begin{enumerate}
\def\labelenumi{(\alph{enumi})}
\item
  若 H 不是闭的, 则 \(\bar{H}\) 是紧的 (《谱理论》第 II 章 §2 引理 1), 因而同构于 \(T^{n}\); 若 G 连通, 则它在 G 中是中心的 (使用《一般拓扑》第 VII 章 §2 命题 5).
\item
  假设 H 是闭的. 令 a 为 G 的一个元素, C(a) 为 G 中与 a 交换的元素的集合. 假设 H ⊕ C(a).
\end{enumerate}

若 H 同构于 R, 则 \(\mathrm{C}(a) \cap \mathrm{H} = \{e\}\). (考虑 H 的自同构 \(\alpha: h \mapsto a^{-1}ha\).) 若 H 同构于 T, 则 \(\mathrm{C}(a) \cap \mathrm{H}\) 有两个元素. (再次考虑 \(\alpha\), 并使用《一般拓扑》第 VII 章 §2 命题 6.) 若 G 连通, 则第二种情形不可能发生 (使用 \((a)\)).

\begin{enumerate}
\def\labelenumi{(\alph{enumi})}
\setcounter{enumi}{2}
\tightlist
\item
  在 (b) 的假设下, 再假设 G/H 是交换的.
\end{enumerate}

于是 \(\mathrm{G} = \mathrm{C}(a)\) . H. (注意从 H 到 H 的映射 \(h \mapsto h^{-1}\alpha(h)\) 是满射.)

\begin{enumerate}
\def\labelenumi{\arabic{enumi}.}
\setcounter{enumi}{21}
\tightlist
\item
  设 G 是实李群, H 是正规的一维积分子群, A 是 G 的闭子群, 且 AH 在 G 中不是闭的. 则 H 是 AH 的单位元连通分支中的中心子群. (将问题归约为 G = AH 且 G 连通的情形. 令 B 为 A 中与 H 交换的元素所成的集合, 则 B 是 G 的正规子群. 对 B 取商, 将问题归约为 B = \(\{e\}\) 的情形. 由于 G 中每个换位子都与 H 交换, A 于是是交换的. 假设 H 是闭的并使用练习 21 (c), 证明 AH 将是闭的, 矛盾. 因而 H 不是闭的, 可以使用练习 21 (a).)
\end{enumerate}

¶ 23. 设 G 是实李群, \(c = (\mathrm{U}, \phi, \mathrm{E})\) 是以单位元 e 为中心的 G 的流形上的一个图. 若 \(e\)\(x \in \mathrm{U}\), 令 \(|x|_c\) 表示 Banach 空间 E 中 \(\phi(x)\) 的范数.

\begin{enumerate}
\def\labelenumi{(\alph{enumi})}
\tightlist
\item
  证明若 \(c' = (\mathbf{U}', \phi', \mathrm{E}')\) 是以 e 为中心的另一个图, 则存在常数 \(e\)\(\lambda > 0\) 和 \(\mu > 0\), 使得
\end{enumerate}

\[
\mu | x | _ {c} \leqslant | x | _ {c ^ {\prime}} \leqslant \lambda | x | _ {c}
\]

对所有充分接近 e 的 \(x \in G\) 都成立.

\begin{enumerate}
\def\labelenumi{(\alph{enumi})}
\setcounter{enumi}{1}
\tightlist
\item
  证明对所有 \(\rho > 0\), 存在包含于  的 e 的邻域 \(\mathbf{U}_{\rho}\), 使得对 \(e\) 中的 \(\mathbf{U}\)\(x, y\), \(\mathbf{U}_{\rho}\)\((x, y) \in \mathbf{U}_{\rho}\), 且
\end{enumerate}

\[
\left| (x, y) \right| _ {c} \leqslant \rho . \inf (\left| x \right| _ {c}, \left| y \right| _ {c}).
\]

\begin{enumerate}
\def\labelenumi{(\alph{enumi})}
\setcounter{enumi}{2}
\tightlist
\item
  假设 \(\mathbf{G}\) 是有限维的. 令 \(\Gamma\) 为 \(\mathbf{G}\) 的离散子群. 对 \(0 < \rho < 1\) 使用 (b), 并选取相对紧的 \(\mathrm{U}_{\rho}\). 集合 \(\mathrm{U}_{\rho} \cap \Gamma\) 是有限的. 令
\end{enumerate}

\[
\{e = \gamma_ {0}, \gamma_ {1}, \dots , \gamma_ {m} \}
\]

是它的元素, 并按如下方式编号:

\[
\left| \gamma_ {0} \right| _ {c} \leqslant \left| \gamma_ {1} \right| _ {c} \leqslant \dots \leqslant \left| \gamma_ {m} \right| _ {c}.
\]

证明若 \(i \leqslant m\) 且 \(j \leqslant m\), 则换位子 \((\gamma_{i}, \gamma_{j})\) 等于某个  的 \(\gamma_{k}\). 推出由 \(k < \inf(i,j)\)\(\mathrm{U}_{\rho} \cap \Gamma\) 生成的子群是类 \(\leqslant m\) 的幂零群.

由上述结果推出存在 e 的一个邻域 \(\mathbf{V}\), 使得对每个 \(e\) 的离散子群 \(\Gamma\), 都存在包含  的 \(\mathbf{G}\) 的幂零积分子群 \(\mathbf{N}\).\(\mathbf{G}\)\(\mathbf{V} \cap \Gamma\)

\begin{enumerate}
\def\labelenumi{(\alph{enumi})}
\setcounter{enumi}{3}
\tightlist
\item
  假设 G 是连通且有限维的, 并包含离散子群的递增序列 \(D_{n}\), 其并在 G 中稠密. 则 G 是幂零的. (利用 (c), 证明存在一个中心的单参数子群 H, 它与充分大的 n 所对应的 \(D_{n}\) 相交于一个不同于 e 的点. 按  归纳, 根据 H 是闭的还是相对紧的区分两种情形 (练习 21 (a)).
\end{enumerate}

\begin{enumerate}
\def\labelenumi{\arabic{enumi}.}
\setcounter{enumi}{23}
\tightlist
\item
  设 \(G, G'\) 是有限维实李群, f 是从 G 到 \(f\)\(G\)\(G'\) 的满射态射, N 是其核, H 是 G 的积分子群, \(N\)\(H\)\(G\)\(H' = f(H)\).
\end{enumerate}

\begin{enumerate}
\def\labelenumi{(\alph{enumi})}
\item
  假设 N 是有限的. \(H'\) 是闭的, 当且仅当 H 是闭的.
\item
  假设 N 是紧的. 若 H 是闭的, 则 \(H'\) 是闭的.
\end{enumerate}

\begin{enumerate}
\def\labelenumi{\arabic{enumi}.}
\setcounter{enumi}{24}
\tightlist
\item
  设 \(G\) 是有限维实或复李群. 令 \(S \subset L(G)\). 假设 L(G) 是由 S 生成的李代数, 且 S 在伸缩变换下不变.\(L(G)\)\(S\)\(S\)
\end{enumerate}

\begin{enumerate}
\def\labelenumi{(\alph{enumi})}
\item
  令 H 为由 exp S 生成的 G 的子群. 证明 H 在 G 中是开的. (令 A 为满足  的 \(x \in L(G)\) 的集合, B 为由 A 生成的 L(G) 的向量子空间. 证明 \(\exp(Kx) \subset H\)\(L(G)\){[}A, S{]} \(\subset A\), 从而 {[}A, A{]} \(\subset B\). 推出 B 是 L(G) 的子代数, 然后使用 § 4 的命题 3.)\(L(G)\)
\item
  再假设
\end{enumerate}

\[
(x \in S \text {   and   } y \in S) \Rightarrow ((\operatorname{Ad} \exp x) (y) \in S).
\]

则由  生成的  的向量子空间 \(\mathbf{V}\) 等于 \(\mathbf{L}(\mathbf{G})\)\(\mathbf{S}\)\(\mathbf{L}(\mathbf{G})\). (利用 \((a)\), 证明 \([\mathbf{L}(\mathbf{G}), \mathbf{V}] \subset \mathbf{V}\).)

\begin{enumerate}
\def\labelenumi{\arabic{enumi}.}
\setcounter{enumi}{25}
\tightlist
\item
  设 G 是 n 维连通紧复李群, X 是有限维连通复解析流形, \((g, x) \mapsto gx\) 是 G 在 X 上的解析左作用法则. 对所有 \(g \in G\), 令 \(\rho(g)\) 为从 X 到 X 的映射 \(x \mapsto gx\). 假设对所有不同于 e 的 \(g \in G\), \(e, \rho(g) \neq Id_x\). 则 G 在 X 中的每个轨道都是 X 的 n 维闭子流形. (令 \(x \in X\), H 为 x 的稳定子, H' 为 H 的单位元连通分支. 对所有 \(H'\)\(g \in H'\), 令 \(u(g)\) 为 \(\rho(g)\) 在 x 处的切映射. 仿照第 3 号命题 6 的论证, 证明对所有 , \(u(g)\) 都是恒等映射. 使用 § 1 练习 4 以及 X 的连通性, 推出 \(g \in H'\)\(H' = \{e\}\).)
\end{enumerate}

¶ 27. 设 G 是紧实李群, M 是  类紧流形, J 是包含 0 的 R 的一个开区间. 令 \((s, \langle x, \xi \rangle) \mapsto (m_{\xi}(s, x), \xi)\) 为从 G × (M × J) 到 M × J 的  类映射, G 通过该映射在 M × J 上左作用.

\begin{enumerate}
\def\labelenumi{(\alph{enumi})}
\item
  令 X 是 M × J 上 \(C^{1}\) 类向量场, 使得对所有 \(M \times J\)\((x, \xi) \in M \times J\), \(\mathbf{X}_{(x, \xi)}\) 的第二投影是 J 的切向量 l. 将 X 用 G 平移并在 G 上积分, 推出存在一个与 X 具有相同性质且进一步在 G 下不变的场 \(X'\).
\item
  证明存在从  到自身的微分同胚 \((x, \xi) \mapsto (h_{\xi}(x), \xi)\), 使得\(\mathbf{M} \times \mathbf{J}\)
\item
  对所有 \(\xi \in \mathbf{J}\), \(h_{\xi}\) 都是从 \(\mathbf{M}\) 到 \(\mathbf{M}\) 的微分同胚;
\end{enumerate}

\begin{enumerate}
\def\labelenumi{(\roman{enumi})}
\setcounter{enumi}{1}
\tightlist
\item
  \(m_{\xi}(s,x) = h_{\xi}(m_0(s,h_{\xi}^{-1}(x)))\) 对 \(s\in \mathbf{G}, x\in \mathbf{M},\xi \in \mathbf{J}\) 成立.
\end{enumerate}

(使用 (a) 和第 8 号定理 5.)

\begin{enumerate}
\def\labelenumi{\arabic{enumi}.}
\setcounter{enumi}{27}
\item
  举出一个有限维实李群 G 以及 G 的两个积分子群 A, B 的例子, 使得 \(\mathrm{L}(\mathrm{G})=\mathrm{L}(\mathrm{A})+\mathrm{L}(\mathrm{B})\), 但 AB 不是 G 的子群 (参见《积分论》第 VII 章 § 3 第 3 号命题 6).
\item
  设 G 是有限维连通复李群, \(G_{0}\) 是底层实李群, H 是 \(G_{0}\) 的积分子群. 证明存在包含 H 的 G 的最小积分子群 \(H^{*}\). 举出 H 在 \(G_{0}\) 中闭而 \(H^{*}\) 在 G 中不闭的例子 (取 \(G = C^{2}/Z^{2}\)).
\item
  设 \(G\) 是紧连通复李群, 因而形如 \(\mathbf{C}^n / D\), 其中 D 是 \(D\)\(\mathbf{C}^n\) 的秩为 \(2n\) 的离散子群.
\end{enumerate}

\begin{enumerate}
\def\labelenumi{(\alph{enumi})}
\item
  证明 G 上每个全纯 1 次微分形式 \(\omega\) 都是不变的. (令 \(\pi: \mathbf{C}^n \to \mathbf{G}\) 为典范态射. 令 \(\zeta_1, \ldots, \zeta_n\) 为 \(\mathbf{C}^n\) 上的坐标函数. 则 \(\pi^*(\omega)\) 形如 \(\sum_j a_j d\zeta_{j}\), 其中 \(a_j\) 是 \(\mathbf{C}^n\) 上在 D 下不变的全纯函数, 因而是常数.)
\item
  令 \(G' = C^{n}/D'\), 其中 \(D'\) 是 \(C^{n}\) 的秩为 2n 的离散子群. 证明每个解析流形同构 \(u: G \to G'\) 都具有形式 \(s \mapsto v(s) + a'\), 其中 v 是李群同构且 \(a' \in G'\). (令 \(\pi': C^{n} \to G'\) 为典范态射. 存在解析流形同构 \(\tilde{u}: C^{n} \to C^{n}\), 使 \(u \circ \pi' \circ \tilde{u}\). 对 G' 上的每个全纯 1 次微分形式 \(\omega'\), \(G'\)\(u^*(\pi^*(\omega'))\) 是 \(C^{n}\) 上的 1-形式, 根据 (a) 在平移下不变. 推出 \(\tilde{u}\) 是仿射映射.)
\end{enumerate}

\exercisesfor{7}

\begin{enumerate}
\def\labelenumi{\arabic{enumi}.}
\item
  设 \(\mathrm{G}\) 是李群, \(\phi: \mathrm{U} \to \mathrm{G}\) 是指数映射, 满足 \(\mathbf{Z}\mathbf{U} \subset \mathbf{U}\), 且对所有 \(\phi(rx) = \phi(x)^r\)\(x \in \mathbf{U}\) 和所有 \(r \in \mathbf{Z}\) 有 . 若 \(p > 0\), 则 \(\phi\) 是从 \(\mathrm{U}\) 到 \(\phi(\mathrm{U})\) 的解析同构. (若 \(\phi(x) = \phi(y)\), 则对所有  有 \(\phi(p^n x) = \phi(p^n y)\), 从而 \(n \in \mathbf{N}\)\(x = y\). 令 \(\mathrm{W}\) 为 \(0\)\(\mathrm{L}(\mathrm{G})\) 中 0 的一个开邻域, 使 \(\phi^{-1}\) 在 \(\phi(\mathrm{W})\) 上解析. 对所有 \(s \in \phi(\mathrm{U})\), 存在 \(n \in \mathbf{N}\) 及  中 s 的一个邻域 \(\mathrm{V}\), 使 \(s\)\(\phi(\mathrm{U})\)\(t \in \mathrm{V} \Rightarrow t^{p^n} \in \phi(\mathrm{W})\).)
\item
  设 G 是李群, \(\phi: \mathrm{U} \to \mathrm{G}\) 是具有命题 3 所述性质的指数映射, V 是 U 中 0 的一个邻域, 满足 \(\mathbf{ZV} \subset \mathrm{V}\), \(\psi: \mathrm{V} \to \mathrm{G}\) 是 \(\phi\) 在 0 处的切映射, 且对所有 \(\psi(nx) = \psi(x)^n\)\(n \in \mathrm{V}\) 和所有 \(n \in \mathbf{Z}\) 有 . 若 \(p > 0\), 则 \(\psi = \phi \mid \mathrm{V}\). (给 L(G) 赋予一个范数. 令 \(x \in \mathrm{V}\). 当 n 趋于  时, \(p^n x\) 趋于 0, 因而存在 \(n\)\(+\infty\)\(\alpha_n > 0\), 使 \(\alpha_n\) 趋于 0 且 \(\|(\phi^{-1} \circ \psi)(p^n x) - p^n x \| \leqslant \alpha_n \|p^n x\|\). 但 \(\psi(p^n x) = \psi(x)p^n\), 所以 \(\|(\phi^{-1} \circ \psi)(x) - x \| \leqslant \alpha_n \|x\|\)。)
\item
  令 U 为 A 的可逆元素的集合, 且 \(U' = 1 + m \subset U\).
\end{enumerate}

\begin{enumerate}
\def\labelenumi{(\alph{enumi})}
\item
  证明 \(\mathbf{U}' \subset \mathbf{U}_f\). (若 \(x = 1 + y\), 其中 \(y \in \mathfrak{m}\), 则由二项式定理, 当 \(x^{p^n}\)\(n\) 趋于 \(+\infty\) 时  趋于 1.)
\item
  证明 \(U_{f}\) 是 U 中像在 A/m 中为单位根的元素的集合. (使用 (a).) 从而在 K 局部紧时恢复 \(U = U_{f}\) 的事实 (此时 A/m 是有限的).
\end{enumerate}

\begin{enumerate}
\def\labelenumi{\arabic{enumi}.}
\setcounter{enumi}{3}
\tightlist
\item
  设 \(n \in \mathbf{N}^*\), p 是素数, G 是 \(p\)\(G\)\(\mathbf{GL}(n, \mathbf{Z}_p)\) 中所有元素在 p \(p\) 时模 p 同余于 1 (而在 \(p \neq 2\)\(p = 2\) 时模 4 同余于 1) 的矩阵所成的集合. 则 G 是 \(G\)\(\mathbf{GL}(n, \mathbf{Z}_p)\) 的开子群.
\end{enumerate}

\begin{enumerate}
\def\labelenumi{(\alph{enumi})}
\item
  证明 G 没有阶为 \(\neq 1\) 的有限阶元素.
\item
  证明 \(\mathbf{GL}(n, \mathbf{Z}_p)\) 的每个有限子群都同构于  时 \(\mathbf{GL}(n, \mathbf{Z} / p\mathbf{Z})\) 的一个子群 (而 \(p \neq 2\) 时同构于 \(\mathbf{GL}(n, \mathbf{Z} / 4\mathbf{Z})\) 的一个子群).\(p = 2\)
\end{enumerate}

¶ 5. (a) 设 \(\Gamma\) 是 \(G = \mathbf{GL}(n, \mathbf{Q}_p)\) 的紧子群. 证明 \(\Gamma\) 的某个共轭子群包含于 \(\mathbf{GL}(n, \mathbf{Z}_p)\). (若 T 是相对于 \(T\) 的 \(\mathbf{Q}_p^n\) 的格 (《交换代数》第 VII 章 § 4 定义 1), 证明 T 在 \(\mathbf{Z}_p\)\(T\)\(\Gamma\) 中的稳定子是 \(\Gamma\) 的开子群, 因而指数有限, 且 \(\sum_{\gamma \in I} \gamma T\) 是一个在 \(\Gamma\) 下不变的格.)

\begin{enumerate}
\def\labelenumi{(\alph{enumi})}
\setcounter{enumi}{1}
\item
  推出 \(\mathbf{G}_f\) 是 \(\mathbf{GL}(n,\mathbf{Z}_p)\) 的各共轭子群的并.
\item
  推出  的每个有限子群 \(\Phi\) 的阶都是 \(\mathbf{GL}(n, \mathbf{Q}_p)\)\(a_n(p)\) 的一个因子, 其中 \(a_n(p)\) 定义为
\end{enumerate}

\[
\begin{array}{l l} a _ {n} (p) = (p ^ {n} - 1) (p ^ {n} - p) \ldots (p ^ {n} - p ^ {n - 1}) & \text {if} p \neq 2 \\ a _ {n} (p) = 2 ^ {n ^ {2}} (2 ^ {n} - 1) (2 ^ {n} - 2) \ldots (2 ^ {n} - 2 ^ {n - 1}) & \text {if} p = 2. \end{array}
\]

(利用 (a), 可假设 \(\Phi \subset \mathbf{GL}(n, \mathbf{Z}_p)\). 然后使用练习 4.)

\begin{enumerate}
\def\labelenumi{(\alph{enumi})}
\setcounter{enumi}{3}
\tightlist
\item
  证明 \(\mathbf{SL}(n, \mathbf{Q}_p)\) 的每个有限子群的阶都是 \(s_n(p)\) 的一个因子, 其中 \(s_n(p) = \frac{a_n(p)}{p - 1}\) (当 \(p \neq 2\)) 且 \(s_n(2) = \frac{a_n(2)}{2}\).
\end{enumerate}

¶ 6. 本练习中, l 表示一个 \(\neq2\) 的素数. 若 a 是一个 \(\neq0\) 的整数, 令 \(v_{l}(a)\) 表示 a 的 l 进赋值, 即满足 \(a\equiv0\pmod{l^{e}}\) 的最大整数 e.

\begin{enumerate}
\def\labelenumi{(\alph{enumi})}
\tightlist
\item
  令 \(m\) 为一个 \(\geqslant 1\) 的整数. 记
\end{enumerate}

\[
\begin{array}{l l} \varepsilon (l, m) = 0 & \text { if } m \neq 0 (\mathrm{mod} (l - 1)) \\ \varepsilon (l, m) = v _ {l} \left(\frac {m}{l - 1}\right) + 1 & \text { if } m \equiv 0 (\mathrm{mod} (l - 1)). \end{array}
\]

证明若 x 是与 l 互素的整数, 则\(x\)\(l\)

\[
v _ {l} (x ^ {m} - 1) \geqslant \varepsilon (l, m)
\]

并且若 x 在循环群 \(x\)\((\mathbf{Z} / l^2\mathbf{Z})^*\) 中的像是该群的生成元, 则等号成立.

\begin{enumerate}
\def\labelenumi{(\alph{enumi})}
\setcounter{enumi}{1}
\tightlist
\item
  令 \(n\) 为一个 \(\geqslant 1\) 的整数. 记
\end{enumerate}

\[
r (l, n) = \sum_ {m = 1} ^ {n} \varepsilon (l, m).
\]

证明

\[
r (l, n) = \left[ \frac {n}{l - 1} \right] + \left[ \frac {n}{l (l - 1)} \right] + \left[ \frac {n}{l ^ {2} (l - 1)} \right] + \dots
\]

其中符号 \([\alpha]\) 表示实数 \(\alpha\) 的整数部分. 证明若 x 是与 l 互素的整数, 则

\[
v _ {l} ((x ^ {n} - 1) (x ^ {n - 1} - 1) \dots (x - 1)) \geqslant r (l, n)
\]

并且若 x 在 \(x\)\((\mathbf{Z} / l^2\mathbf{Z})^*\) 中的像是该群的生成元, 则等号成立.

\begin{enumerate}
\def\labelenumi{(\alph{enumi})}
\setcounter{enumi}{2}
\item
在练习 5 (c) 的记号下, 证明 \(l^{r(l,n)}\) 是当 p 为 \(l\) 的素数时整除所有 \(a_{n}(p)\) 的 l 的最大幂 (或者对所有充分大的 p, 这两种说法等价). (对 \(p\)\(\neq 1\)\(p\)\(x = p\) 应用 (b); 然后选取 p, 使它在 \(p\)\((\mathbf{Z}/l^2\mathbf{Z})^*\) 中的像是该群的生成元; 根据算术级数定理这是可能的.†)
\item
  令 \(\bar{\Gamma}\) 为 \(\mathbf{GL}(n, \mathbf{Q})\) 的有限子群, 令 \(l^e\) 为整除 \(l\)\(\Gamma\) 的阶的最大 l 幂. 证明 \(e \leqslant r(l, n)\). (利用练习 5 证明 \(l^e\) 整除所有 \(a_n(p)\), 然后应用上面的 (c).)
\item
  反之, 证明存在 \(l\) 的有限 l-子群 \(\Gamma_{i,n}\), 其阶为 \(\mathbf{GL}(n,\mathbf{Q})\)\(l^{d,l,n}\). (将其归约为 n 形如 \(n\)\(l^a (l - 1)\) 且 \(a\geqslant 0\) 的情形. 将 \(\mathbf{Q}^{\mathrm{n}}\) 分解为 \(l^a\) 个 \(\mathbf{Q}^{j - 1}\) 的副本之和, 并利用这一分解在 \(\mathbf{Q}^{\mathrm{n}}\) 上定义对称群 \(\mathrm{H}_{l,n}\)\(\mathfrak{S}_{la}\) 与群 \((\mathbf{Z} / l\mathbf{Z})^{a}\) 的半直积  的作用. 令 \(\Gamma_{l,n}\) 为 \(l\)\(\mathrm{H}_{l,n}\) 的一个 Sylow l-子群.)
\end{enumerate}

证明若 n 为偶数, 则 \(n\)\(\Gamma_{l,n}\) 包含于 \(\mathbf{Sp}(n,\mathbf{Z})\) 的某个共轭子群中.

\begin{enumerate}
\def\labelenumi{(\alph{enumi})}
\setcounter{enumi}{5}
\tightlist
\item
  *令 \(\Gamma\) 为 \(\mathbf{GL}(n, \mathbf{Q})\) 的有限子群, 且是一个 l-群. 证明 \(l\)\(\Gamma\) 共轭于 \(\Gamma_{l,n}\) 的某个子群. (首先利用适当的模 p 约化, 证明 \(p\)\(\Gamma\) 的约化是 \(\Gamma_{l,n}\) 的某个共轭子群的约化的有限子群; 然后使用 \(\Gamma\) 在 \(\mathbf{Q}^{n}\) 上的表示的特征标.) 特别地, \(\mathbf{GL}(n, \mathbf{Q})\) 中阶为 \(l^{(l,n)}\) 的每个子群都共轭于 \(\Gamma_{l,n*}\).
\end{enumerate}

¶ 7. 在本练习中, 令 \(v_{2}(a)\) 表示整数 a 的 2 进赋值.\(a\)

\begin{enumerate}
\def\labelenumi{(\alph{enumi})}
\item
  令 \(\Gamma\) 为 \(\mathbf{GL}(n,\mathbf{Q})\) 的有限子群. 证明存在一个系数为整数且正定的二次型, 在 \(\mathbf{Z}\)\(\Gamma\) 下不变. 推出 (用与练习 4 和 5 相同的论证), 对所有充分大的 p, \(p\)\(\Gamma\) 同构于域  上正交群 \(\mathbf{O}(n)\) 的一个子群 (若 \(\mathbf{F}_p\)\(\mathbf{SO}(n)\)\(\Gamma\) 包含于 \(\mathbf{SL}(n,\mathbf{Q})\), 则同构于  的一个子群).
\item
  假设 \(\Gamma\) 包含于 \(\mathbf{SL}(n, \mathbf{Q})\), 令 \(2^e\) 表示整除 \(\Gamma\) 的阶的最大 2 幂. 证明若 n 为奇数, 则 \(n\)\(2^e\) 整除所有整数
\end{enumerate}

\[
b _ {n} (p) = \left(p ^ {n - 1} - 1\right) \left(p ^ {n - 3} - 1\right) \dots \left(p ^ {2} - 1\right)
\]

(使用 (a) 和《代数》第 IX 章 § 6 的练习 13.)\(p\)

证明若 n 为偶数且 p 是充分大的素数, 则 \(n\)\(p\)\(2^e\) 整除整数\(b_n(p)\)

\[
\begin{array}{l} (p ^ {n} - 1) (p ^ {n - 2} - 1) \dots (p ^ {2} - 1) / (p ^ {n / 2} + 1) \\ (p ^ {n} - 1) (p ^ {n - 2} - 1) \dots (p ^ {2} - 1) / (p ^ {n / 2} - 1). \end{array}
\]

(与 n 为奇数时相同的方法.)\(n\)

\begin{enumerate}
\def\labelenumi{(\alph{enumi})}
\setcounter{enumi}{2}
\tightlist
\item
在 (b) 的假设下, 记
\end{enumerate}

\[
r (2, n) = n + \left[ \frac {n}{2} \right] + \left[ \frac {n}{4} \right] + \dots = n + v _ {2} (n!).
\]

令 A 为一个 \(\geqslant 3\) 的整数. 证明 \(2^{r(2,n)-1}\) 是整除所有  时的 \(b_n(p)\) 的 2 的最大幂. (与练习 6 中相同的方法†; 使用存在 \(p \geqslant A\)\(p \geqslant A\) 且 \(p \equiv 5 \pmod{8}\) 的素数这一事实.) 推出不等式 \(e \leqslant r(2,n) - 1\).

\begin{enumerate}
\def\labelenumi{(\alph{enumi})}
\setcounter{enumi}{3}
\tightlist
\item
反之, 令 \(C_{n}\) 为由置换矩阵和系数为  的对角矩阵生成的 \(\mathbf{GL}(n,\mathbf{Z})\) 的子群. \(\pm1\)\(C_{n}\) 的阶为 \(2^{n}n!\), 且 \(v_{2}(2^{n}n!) = r(2,n)\). 若 \(\Gamma_{2,n}\) 表示 \(C_{n}\) 的一个 Sylow 2-子群与 \(\mathbf{SL}(n,\mathbf{Z})\) 的交, 推出 \(\Gamma_{2,n}\) 的阶为 \(2^{n(2,n)-1}\).
\end{enumerate}

\begin{enumerate}
\def\labelenumi{\arabic{enumi}.}
\setcounter{enumi}{7}
\tightlist
\item
  令 \(n\) 为一个 \(\geqslant 1\) 的整数. 记
\end{enumerate}

\[
\mathrm{M} (n) = \prod_ {l} 1 ^ {r (1, n)},
\]

其中乘积遍历所有素数 l, 而 \(l\)\(r(l, n)\) 如练习 6 和 7 中定义.

于是 \(\mathbf{M}(1) = 2\), \(\mathbf{M}(2) = 2^3.3 = 24\), \(\mathbf{M}(3) = 2^4.3 = 48\),

\[
\mathrm{M} (4) = 2 ^ {7}. 3 ^ {2}. 5 = 5 7 6 0.
\]

由练习 6 和 7 推出, \(\mathbf{GL}(n,\mathbf{Q})\) (或 \(\mathbf{GL}(n,\mathbf{Z})\), 这两者没有区别) 的有限子群的阶的最小公倍数等于 \(\mathrm{M}(n)\). 对 \(\mathbf{SL}(n,\mathbf{Q})\) 考虑同一问题, 其中以  代替 \(\mathrm{M}(n)\).\(\frac{1}{2}\mathrm{M}(n)\)

\begin{enumerate}
\def\labelenumi{\arabic{enumi}.}
\setcounter{enumi}{8}
\item
  假设 \(\mathbf{K}\) 是局部紧的. 令 \(\mathrm{G} = \mathbf{GL}(n,\mathbf{K})\). 令 \(\mathbf{G}_1\) 为保持  的某个相对于 A 的格不变的 \(g\in \mathbf{G}\) 的集合. 令 \(\mathbf{K}^n\)\(\mathbf{G}_2\) 为在 G 中生成相对紧子群的 \(g\in \mathbf{G}\) 的集合. 令 \(\mathbf{G}_3\) 为其在 \(g\in \mathbf{G}\)\(\mathbf{K}\) 的代数闭包中的特征值绝对值为 1 的  的集合. 则 \(\mathbf{G}_1 = \mathbf{G}_2 = \mathbf{G}_3 = \mathbf{G}_f\). (使用练习 5 (a) 的论证.)
\item
  假设 K 是局部紧的. 令 G 为 K 上 n 维标准群. 令 \(\mu\) 为加法群 \(K^{n}\) 上的 Haar 测度. 证明 \(\mu|G\) 是 G 上的左、右 Haar 测度. (使用 G 是各 \(G(\alpha_{\lambda})\) 的逆极限这一事实以及《积分论》第 VII 章 §1 命题 7.)
\end{enumerate}

\exercisesfor{8}

\begin{enumerate}
\def\labelenumi{\arabic{enumi}.}
\item
  从  到  的映射 \(z \mapsto \bar{z}\) 是复李群 \(\mathbf{C}\) 的非解析连续自同构.\(\mathbf{C}\)\(\mathbf{C}\)
\item
  令 G 为由实数序列 \((\lambda_{1}, \lambda_{2}, \ldots)\) 组成的 Hilbert 空间, 满足 \(\sum_{i} \lambda_{i}^{2} < +\infty\). 将 G 视为实李群. 令 \(G_{n}\) 为由 \((\lambda_{1}, \lambda_{2}, \ldots) \in G\) 组成的集合, 其中对  有 \(\lambda_{m} \in \frac{1}{m} Z\). 各 \(1 \leqslant m \leqslant n\)\(G_{n}\) 是 G 的闭李子群, 因而 \(H = \bigcap_{n} G_{n}\) 是 G 的闭子群. 但此群全不连通且非离散, 因而不是 G 的李子群.
\end{enumerate}

¶ 3. 在 \(Q_{p} \times Q_{p}\) 中, 每个闭集都可以由一族解析方程定义. 推出定理 2 的推论 2 (ii) 在省去 I 有限的假设时不再成立.

¶ 4. 设 G 是有限维实李群, A 是 G 的一个子群. 称 L(G) 中的元素 x 是 A-可达的, 如果对 e 的每个邻域 U, 都存在从 \(\alpha\){[}0, 1{]} 到 A 的连续映射 , 使 \(\alpha(0) = e\), 且对  有 \(\alpha(t) \in \exp(tx)\) . U. 令 b 为 L(G) 中 A-可达元素的集合.\(0 \leqslant t \leqslant 1\)

\begin{enumerate}
\def\labelenumi{(\alph{enumi})}
\item
  证明 \(\mathfrak{h}\) 是 \(\mathbf{L}(\mathbf{G})\) 的李子代数. (使用 §6 的练习 19.)
\item
  令 H 是 G 的积分子群, 满足 \(\mathrm{L}(\mathrm{H}) = \mathfrak{h}\) . 证明 \(H \subset A\) . (令 I = \(\{-1, 1\}\), 令 \(R'\) 赋予欧氏范数. 令 \((x_{1}, \ldots, x_{r})\) 为 h 的一个基. 构造从 I 到 A 的连续映射 \(\alpha_{1}, \ldots, \alpha_{r}\) 以及从  到 R 的连续映射 \(f_{1}, \ldots, f_{r}\), 使得对 \(I'\) 中的 \(t = (t_{1}, \ldots, t_{r}) \in I'\) 有
\end{enumerate}

\[
\begin{array}{c} \alpha_ {1} (t _ {1}) \ldots \alpha_ {r} (t _ {r}) = \exp (f _ {1} (t) x _ {1}) \ldots \exp (f _ {r} (t) x _ {r}) \\ \| t - (f _ {1} (t), \ldots , f _ {r} (t)) \| \leqslant \frac {1}{2}. \end{array}
\]

然后应用以下定理: 令 \(f\) 为从 \(\mathbf{I}^r\) 到 \(\mathbf{R}^r\) 的连续映射, 使对所有  有 \(\| f(x) - x \| \leqslant \frac{1}{2}\); 则 \(x \in \mathbf{I}^r\)\(f(\mathbf{I}^r)\) 包含 \(\mathbf{R}^r. \dagger\) 中 0 的一个邻域

\begin{enumerate}
\def\labelenumi{(\alph{enumi})}
\setcounter{enumi}{2}
\item
  推出 \(\mathfrak{h}\) 是 A 在 e 处相切的子代数.\(e\)
\item
  证明若 A 是弧连通的, 则 A = H.†
\end{enumerate}

¶ 5. 设 G 是 Hausdorff 拓扑群, H 是 G 的闭子群, \(\pi\) 是从 G 到 G/H 的典范映射. 假设 H 是有限维实李群. 则 G/H 中 \(\pi(e)\) 的某个邻域 U 上存在连续映射 \(\sigma\) 到 G, 使 \(\pi \circ \sigma = \mathrm{Id}_{\mathrm{U}}\). (令 \(\rho\) 为 H 在 \(\mathbf{GL}(n,\mathbf{R})\) 上的解析线性表示, 它是局部同胚 (§ 6, 定理 1 的推论). 令 f 为 G 上满足 \(f\)\(\geqslant 0\) 的连续函数, 在 e 处等于 1, 并在 e 的某个充分小邻域外为零. 令 \(e\)\(e\)\(ds\) 为 H 的左 Haar 测度. 对 \(x \in G\), 记

\[
g (x) = \int f (x s) \rho (s) ^ {- 1} d s \in \mathbf {M} _ {n} (\mathbf {R}).
\]

于是对 \(g(xt) = g(x)\varphi(t)\)\(x \in G\) 和 \(t \in H\) 有 . 若 V 足够小,

\[
g (x) \in \mathbf {G L} (n, \mathbf {R})
\]

当 x 充分接近 e 时, \(x\). 最后使用如下事实: 要证明的定理对 \(e\)\(\mathbf{GL}(n, \mathbf{R})\) 和 \(\varphi(\mathbf{H})\) 在局部成立.

\begin{enumerate}
\def\labelenumi{\arabic{enumi}.}
\setcounter{enumi}{5}
\item
  设 G 是一个 C＾r 类群 (§ 5 练习 1), 其中 \(C^{r}\)\(r \geqslant 2\). 在 G 上存在唯一的实李群结构 S, 使 S 的 C＾r 类底层流形结构就是给定结构. (S 的唯一性由定理 1 的推论 1 得出. 令 L(G) 为 § 5 练习 2 中与 G 相关的可赋范李代数. 存在一个实李群芽 \(C^{r}\)\(L(G)\)\(G^{r}\) 以及从 \(L(G^{\prime})\) 到 L(G) 的同构 h (§ 4 定理 3). 仿照 § 4 第 1 号验证: 存在 \(L(G)\) 中  的对称开邻域 \(G^{\prime\prime}\), 以及从 \(e_{G}\) 到 G 的 C＾r 类映射 \(G^{r}\)\(\phi\), 满足 \(C^{r}\), 且对 \(G^{\prime\prime}\) 中的 \(T_{e}(\phi) = h\) 有 \(\phi(g_{1}g_{2}) = \phi(g_{1})\phi(g_{2})\). 缩小 \(g_{1}, g_{2}\)\(G^{\prime\prime}\)\(G^{\prime\prime}\) 后, 可假设 \(V = \phi(G^{r})\) 在 G 中开, 且 \(\phi\) 是从流形 \(C^{r}\)\(G^{\prime\prime}\) 到流形 V 的 C＾r 类同构. 因而 V 上存在一个实李群芽结构, 使其 C＾r 类底层流形结构为给定结构. 对所有 \(C^{r}\)\(g \in G\), Int g 定义了从 \(V \cap (g^{-1}Vg)\) 到 \((gVg^{-1}) \cap V\) 的解析映射 (定理 1). 根据 § 1 的命题 18, G 上存在实李群结构 S, 在 e 的某个开邻域上与 V 诱导相同的解析结构. 通过平移, S 在 G 上的 C＾r 类底层流形结构就是给定结构.)\(C^{r}\)
\item
  设 \(G\) 是实李群, H 是 G 的闭子群.\(H\)\(G\)
\end{enumerate}

\begin{enumerate}
\def\labelenumi{(\alph{enumi})}
\tightlist
\item
  令 \(\mathfrak{h}\) 为所有满足对  都有  的 \(x\in \dot{\mathbf{L}} (\mathbf{G})\) 所成的集合. 则 \(\exp (tx)\in \dot{\mathbf{H}}\)\(t\in \mathbf{R}\)\(\mathfrak{h}\) 是 \(\mathrm{L}(\mathrm{G})\) 的李子代数. (使用 §6 的命题 8.)
\end{enumerate}

\begin{enumerate}
\def\labelenumi{(\alph{enumi})}
\setcounter{enumi}{1}
\tightlist
\item
  假设 \(\mathbf{H}\) 是局部紧的. 证明 \(\mathbf{H}\) 是 \(\mathbf{G}\) 的李子群. (首先证明 \(\mathfrak{h}\) 是有限维的, 方法是证明 \(\mathfrak{h}\) 中存在一个相对紧邻域. 然后仿照定理 2 的证明, 选取 \(\mathrm{V}_2\), 使 \((\exp \mathrm{V}_2) \cap \mathrm{H}\) 相对紧.)
\end{enumerate}

\exercisesfor{9}

\begin{enumerate}
\def\labelenumi{\arabic{enumi}.}
\item
  令 \(\mathrm{G} = \mathbf{SO}(3, \mathbf{R})\). 令 \(\Delta_{1}, \Delta_{2}\) 为 \(R^{3}\) 中的两条正交直线, 令 \(H_{i}\) 为 G 中绕 \(\Delta_{i} (i = 1, 2)\) 的旋转所成的子群 . 则 \([\mathrm{L}(\mathrm{H}_{1}), \mathrm{L}(\mathrm{H}_{2})]\) 是 G 的一个 1 维李子代数, 不同于 \((\mathrm{H}_{1}, \mathrm{H}_{2})\) 在 e 处相切的李子代数.
\item
  假设 K 是超度量域. 设 G 是有限维李群, g 是其李代数, A 是 G 的有限子集. 则 \(Z_{G}(A)\) 是 G 的李子群, 其李代数为 \(\mathfrak{z}_{\mathfrak{g}}(A)\). (仿照命题 8.)
\item
  设 G 是连通实或复李群. G 的中心 Z 是 G 的李拟子群, 且 L(Z) 是 L(G) 的中心.
\item
  假设 K 是超度量域且 p \textgreater{} 0 (采用 §7 的记号). 设 G 是有限维李群, A 是 G 的自同构群, B 是相应的 L(G) 的自同构群. 令 \(\mathrm{G}^{\mathbf{A}}\) (分别为 \(\mathrm{L}(\mathrm{G})^{\mathbf{B}}\)) 表示 G (分别为 L(G)) 中在 A (分别为 B) 下不动的元素的集合. 则 \(G^{A}\) 是 G 的李子群, 其李代数为 \(\mathrm{L}(G)^{\mathbf{B}}\). (使用对数映射.)
\item
  设 \(G\) 是有限维连通实或复李群. 令 \((G_0, G_1, \ldots)\) 为 (第 II 章 § 4 练习 18 的) 上中心列, 令 \((g_0, g_1, \ldots)\) 为 L(G) 的上中心列 (第 I 章 § 1 第 6 号). 则对所有 \(L(G)\)\(i\), \(G_i\) 是 G 的李子群, 且 \(G\)\(L(G_i) = g_i\).
\item
  令 \(\Gamma\) 为 §4 练习 5 (b) 中定义的 3 维单连通幂零实李群. 令 \(\alpha \in \mathbf{R}\) 为无理数. 令 P 为 \(P\)\(\Gamma \times \mathbf{R}\) 的离散子群, 由 \(((0,0,x),\alpha x)\) (\(x \in \mathbf{Z}\)) 组成. 令 \(G = (\Gamma \times \mathbf{R}) / P\). 证明 \((G,G)\) 在 G 中不是闭的.\(G\)
\item
  \begin{enumerate}
  \def\labelenumii{(\alph{enumii})}
  \tightlist
  \item
    令 \(G\) 为连通半单实李群, Z 为其中心, \(Z\)\(\rho\) 为 G 在有限维复向量空间上的连续线性表示. 存在整数 p, 使得对所有 \(G\)\(p\)\(z \in Z\), \(\rho(z)\) 都可对角化, 且 \(\rho(z)\) 的所有特征值都是 p 次单位根. (可假设 \(p\)\(\rho\) 不可约. 根据 Schur 引理, \(\rho(z)\) 是标量. 另一方面, 对所有  有 det \(\rho(g) = 1\), 因为 \(g \in G\)\(G = \mathcal{D}G\).)
  \end{enumerate}
\end{enumerate}

\begin{enumerate}
\def\labelenumi{(\alph{enumi})}
\setcounter{enumi}{1}
\tightlist
\item
  推出若 G 具有忠实的有限维连续线性表示, 则 Z 是有限的.\(G\)\(Z\)
\end{enumerate}

¶ 8. 设 G 是有限维连通实李群, g 是其李代数, n 是 g 的最大幂零理想, 且 h = {[}g, g{]} + n.

\begin{enumerate}
\def\labelenumi{(\alph{enumi})}
\item
  \(\mathfrak{h}\) 是 \(\mathfrak{g}\) 的特征理想; \(\mathfrak{h}\) 的根为 n; 对所有 \(n\)\(x \in \mathfrak{h}\), \(\operatorname{Tr} \operatorname{ad}_{\mathfrak{h}} x = 0\); 对于 \(l\)\(\mathfrak{g}\) 的每个 Levi 截面 l, \(\mathfrak{h} = l + n\).
\item
  假设 G 是单连通的. 令 Z 为其中心, H 为李代数为 b 的 G 的李子群, \(\phi\) 为从 G 到 G/H 的典范态射. 则 \(\phi(Z)\) 在 G/H 中是离散的. (群 G 是半单李群 S 与其根 R 的半直积. 令 \(z \in Z\). 则 \(z = y^{-1}x\), 其中 \(x \in R\), y 属于 S 的中心. 存在整数 p, 使 \(Ad_{q}y\) 的特征值, 从而 \(Ad_{q}x\) 的特征值也是 p 次单位根 (练习 7). 推出若 N 是李代数为 n 的 G 的李子群, 则 R 中存在 e 的一个邻域 U, 满足 U = UN = NU,
\end{enumerate}

\[
\mathrm{Z} \cap (\mathrm{SU}) \subset \mathrm{SN} = \mathrm{H}.
\]

\begin{enumerate}
\def\labelenumi{(\alph{enumi})}
\setcounter{enumi}{2}
\item
  现在不再假设 G 是单连通的. 令 H 为李代数为 b 的 G 的积分子群. 证明 H 是 G 的幺模正规李子群, 具有幂零根, 且 G/H 是交换的. (使用 (a) 和 (b).)
\item
  推出若  在 G 中稠密, 则 G 的根是幂零的.
\end{enumerate}

\begin{enumerate}
\def\labelenumi{\arabic{enumi}.}
\setcounter{enumi}{8}
\tightlist
\item
  令 \(G\) 为 \(\mathbf{SL}(2, \mathbf{R})\) 的万有覆盖. 将典范态射 \(G \to \mathbf{SL}(2, \mathbf{R})\) 的核与 \(\mathbf{Z}\) 识别 (§ 6, 练习 2). 令 a 为 \(a\)\(\mathbf{T}^n\) 中幂在 \(\mathbf{T}^n\) 中处处稠密的元素 (《一般拓扑》第 VII 章 §1 命题 7 的推论 2). 令 D 为 \(D\)\(G \times \mathbf{T}^n\) 中由 (1, a) 生成的离散子群. 令 \(H = (G \times \mathbf{T}^n)/D\). 则
\end{enumerate}

\[
\mathrm{L} (\mathrm{H}) = \mathfrak {s l} (2, \mathbf {R}) \times \mathbf {R} ^ {n}
\]

且李代数为  的 H 的积分子群 \(\mathbf{H}'\) 同构于 \(\mathbf{H}\)\(\mathfrak{sl}(2,\mathbf{R})\)\(\mathbf{G}\), 并在 \(\mathbf{H}\) 中稠密. \(\mathrm{D}^n\mathrm{H}' = \mathrm{H}'\) 对所有 \(n\geqslant 0\) 成立.

¶ 10. 设 G 是有限维实或复李群. 令

\[
p = \dim [ L (G), L (G) ].
\]

令 \((\mathrm{G},\mathrm{G})\) 具有其作为 \(\mathbf{G}\) 的积分子群的结构. 存在 \(\mathbf{V}\)\(e\)\((\mathrm{G},\mathrm{G})\) 中 e 的一个邻域 \(\mathbf{V}\), 使得 \(p\) 中的每个元素都是 \(\mathbf{G}\) 中元素的 p 个换位子的乘积. (令 \(x_{1},y_{1},\ldots ,x_{p},y_{p}\) 为 \(\mathbf{L}(\mathbf{G})\) 中的元素, 使 \([x_1,y_1]\) 构成 \(\mathbf{L}(\mathbf{G})\) 的一个基. 对于  中的 \(s,t\), 记\(\mathbf{R}\)

\[
\rho_ {i} (s, t) = (\exp s y _ {i}) ^ {- 1} (\exp t x _ {i}) ^ {- 1} (\exp s y _ {i}) (\exp t x _ {i}).
\]

将隐函数定理应用于映射

\[
\left(s _ {1}, t _ {1}, \dots , s _ {p}, t _ {p}\right) \mapsto \rho_ {1} \left(s _ {1}, t _ {1}\right) \dots \rho_ {p} \left(s _ {p}, t _ {p}\right)
\]

从 \(\mathbf{R}^{2p}\) 到 (G, G).

\begin{enumerate}
\def\labelenumi{\arabic{enumi}.}
\setcounter{enumi}{10}
\item
  设 G 是 B = Z/2Z 与 K 按李群 K 的自同构 \(x \mapsto -x\) 对应的半直积. 则 L(K) 是 L(G) 的理想, L(B) = \{0\}, 但 (K, B) = K.
\item
  设 \((e_1, e_2, e_3)\) 是 \(\mathbf{K}^3\) 的典范基. 考虑 \(\mathbf{K}^3\) 上的幂零李代数结构, 使得 \([e_1, e_2] = e_3\), \([e_1, e_3] = [e_2, e_3] = 0\). 赋予 \(\mathbf{K}^3\) 相应的群法则 (第 5 号),
\end{enumerate}

\[
(x, y, z) (x ^ {\prime}, y ^ {\prime}, z ^ {\prime}) = (x + x ^ {\prime}, y + y ^ {\prime}, z + z ^ {\prime} + \frac {1}{2} (x y ^ {\prime} - y x ^ {\prime})).
\]

映射

\[
(\lambda , \mu , \nu) \mapsto (\exp \lambda e _ {1}) (\exp \mu e _ {2}) (\exp \nu (e _ {1} + e _ {3}))
\]

从 \(K^{3}\) 到 G 的映射既不是满射 (证明 \((0,1,1)\) 不在其像中), 也不是单射 (证明 \((0,1,0)\) 与 \((1,1,-1)\) 的像相同).

\begin{enumerate}
\def\labelenumi{\arabic{enumi}.}
\setcounter{enumi}{12}
\tightlist
\item
  \begin{enumerate}
  \def\labelenumii{(\alph{enumii})}
  \tightlist
  \item
    在 \(\mathbf{R}^3\) 上定义如下乘法:
  \end{enumerate}
\end{enumerate}

\[
(x, y, z). (x ^ {\prime}, y ^ {\prime}, z ^ {\prime}) = (x + x ^ {\prime} \cos z - y ^ {\prime} \sin z, y + x ^ {\prime} \sin z + y ^ {\prime} \cos z, z + z ^ {\prime}).
\]

证明由此得到一个可解实李群 \(\mathbf{G}\), 使得 \(\mathrm{DG} = \mathbf{R}^2\times \{0\}\).  的中心 \(\mathbf{Z}\) 为 \(\mathbf{G}\)\(\{0\} \times \{0\} \times 2\pi \mathbf{Z}\).

\begin{enumerate}
\def\labelenumi{(\alph{enumi})}
\setcounter{enumi}{1}
\tightlist
\item
  对于 \((x,y,z)\in G\), 令 \(\pi (x,y,z)\) 为映射
\end{enumerate}

\[
(\lambda , \mu) \mapsto (\lambda \cos z - \mu \sin z + x, \lambda \sin z + \mu \cos z + y)
\]

从 \(\mathbf{R}^2\) 到 \(\mathbf{R}^2\) 的映射. 证明 \(\pi\) 是从 G 到仿射群 \(G\) 的一个李子群 \(\mathbf{G}'\) 的态射, 该子群由 \(\mathbf{R}^2\) 的平移和旋转生成. 证明 \(\mathbf{R}^2\)\(\operatorname{Ker} \pi = \mathbf{Z}\).

\begin{enumerate}
\def\labelenumi{(\alph{enumi})}
\setcounter{enumi}{2}
\item
  我们典范地将 \(\mathrm{L}(\mathrm{G}) = \mathrm{T}_{(0,0,0)}\mathrm{G}\) 与 \(\mathbf{R}^3\) 识别. 令 \((e_1, e_2, e_3)\) 为 \(\mathbf{R}^3\) 的典范基. 证明 \([e_1, e_2] = 0\), \([e_3, e_1] = e_2\), \([e_3, e_2] = -e_1\).
\item
  证明当 \(c \neq 0\) 时,
\end{enumerate}

\[
\exp_ {\mathrm{G}} (a, b, c) = \left(\frac {1}{c} (a \sin c + b \cos c - b), \frac {1}{c} (- a \cos c + b \sin c + a), c\right).
\]

推出对所有 \(u \in L(G)\), 都有 \(Z \subset \exp(\mathbf{R}u)\). 证明 \(\exp_G: L(G) \to G\) 既不是单射也不是满射.

\begin{enumerate}
\def\labelenumi{\arabic{enumi}.}
\setcounter{enumi}{13}
\tightlist
\item
  \begin{enumerate}
  \def\labelenumii{(\alph{enumii})}
  \tightlist
  \item
    令 \(\mathbf{G} = \mathbf{GL}(n, \mathbf{C})\). 证明 \(\exp_{\mathbf{G}}\) 是满射. (使用《谱理论》第 I 章 § 4 第 8 号的全纯函数演算.)
  \end{enumerate}
\end{enumerate}

\begin{enumerate}
\def\labelenumi{(\alph{enumi})}
\setcounter{enumi}{1}
\tightlist
\item
  令 \(\mathbf{G}' = \mathbf{SL}(2, \mathbf{C})\). 证明若 \(\sigma \in \mathbf{C}^*\), 则元素 \(\left( \begin{array}{cc} -1 & \sigma \\ 0 & -1 \end{array} \right)\)
\end{enumerate}

\(G'\) 不属于 \(\exp_{G'}\) 的像.

\begin{enumerate}
\def\labelenumi{\arabic{enumi}.}
\setcounter{enumi}{14}
\tightlist
\item
  \begin{enumerate}
  \def\labelenumii{(\alph{enumii})}
  \tightlist
  \item
    令 \(x \in \mathfrak{sl}(2, \mathbf{R})\), \(\Delta = \det x\). 证明
  \end{enumerate}
\end{enumerate}

\[
e ^ {x} = \operatorname{ch} \sqrt {- \Delta}. I + \frac {\operatorname{sh} \sqrt {- \Delta}}{\sqrt {- \Delta}} \cdot x \quad \mathrm{if} \Delta <   0
\]

\[
e ^ {x} = \cos {\sqrt {\Delta}}. I + \frac {\sin {\sqrt {\Delta}}}{\sqrt {\Delta}}. x \qquad \mathrm{if} \Delta > 0.
\]

\begin{enumerate}
\def\labelenumi{(\alph{enumi})}
\setcounter{enumi}{1}
\tightlist
\item
  令 \(\mathbf{H} = \mathbf{SL}(2, \mathbf{R})\). 证明 \(g = \begin{pmatrix} \lambda & 0 \\ 0 & \lambda^{-1} \end{pmatrix} \in \mathbf{H}\) 是 \(\exp_{\mathbf{H}}\) 的像, 当且仅当 \(\lambda > 0\) 或 \(\lambda = -1\). 若 \(\lambda = 0\), 则所有包含 g 的单参数子群都相等.\(g\)
\end{enumerate}

\begin{enumerate}
\def\labelenumi{\arabic{enumi}.}
\setcounter{enumi}{15}
\tightlist
\item
  设 G 是有限维李群. 证明存在 L(G) 的一个基 \((x_{1},\ldots,x_{n})\), 使对所有 i, \(\exp(\mathbf{R}x_{i})\) 都是 G 的李子群. (使用《谱理论》第 II 章 §2 引理 1.)
\end{enumerate}

¶ 17. (a) 令 c 表示具有基 \((a, b, c)\) 的实可解李代数, 满足 \([a, b] = c\), \([a, c] = -b\), \([b, c] = 0\). 令 d 表示具有基 \((a, b, c, d)\) 的实可解李代数, 满足 \([a, b] = c\), \([a, c] = -b\), \([b, c] = d\), \([a, d] = [b, d] = [c, d] = 0\). 令 g 为实可解李代数. 若 g 含有非零元素 x, y, z, 满足 \([x, y] = z\) 且 \([x, z] = -y\), 则 g 含有同构于 c 或 d 的子代数.~(令 \(a_{1} = y\), \(b_{1} = z\), \(c_{1} = [y, z]\); 递归地由 \(a_{i}\)\(b_{i}\)\(c_{i}\)\(a_{i} = [a_{i-1}, c_{i-1}]\), \(b_{i} = [b_{i-1}, c_{i-1}]\), \(c_{i} = [a_{i}, b_{i}]\) 定义 . 考虑满足 \(c_{k} = 0\) 的最小整数 k.)

\begin{enumerate}
\def\labelenumi{(\alph{enumi})}
\setcounter{enumi}{1}
\item
  设 g 和 \(g^{*}\) 是实可解李代数, \(\phi\) 是从 g 到 \(g^{*}\) 的满射同态. 若 \(g^{*}\) 含有同构于 c 或 d 的子代数, 则 g 也具有同样性质.
\item
  设 \(\mathfrak{h}\) 是复可解李代数. 考虑 \(\mathfrak{h}\) 的伴随表示的 Jordan-Hölder 列; 该列的商定义了 \(\mathfrak{h}\) 的一维表示, 从而定义了 \(\mathfrak{h}\) 上的线性形式. 这些只依赖于 \(\mathfrak{h}\) 的线性形式称为 \(\mathfrak{h}\) 的根. 若 \(\mathfrak{h}'\) 是实可解李代数, 则 \(\mathfrak{h}'\) 的根是  的根在 \(\mathfrak{h}'\) 上的限制.\(\mathfrak{h}' \otimes_{\mathbf{R}} \mathbf{C}\)
\end{enumerate}

然后令 G 为有限维单连通可解实李群, g 为其李代数. 证明以下条件等价: (α) exp₄ 是单射; (β) exp₅ 是满射; (γ) exp₆ 是双射; (δ) exp₇ 是从解析流形 L(G) 到解析流形 G 的同构; (ε) L(G) 不含同构于 c 或 d 的子代数; (ζ) L(G) 没有商代数含有同构于 c 的子代数; (η) g 的每个根都具有 \(\phi + i\delta'\) 的形式, 其中 \(\phi\), \(\phi'\) 属于 \(g^*\), 且 \(\phi'\) 与 \(\phi\) 成比例; (θ) 对所有 \(x \in g\),  的唯一纯虚特征值 (在 C 中) 是 0.†
